\documentclass[10pt,a4paper]{article}
\usepackage{amsmath,amssymb,amsthm,graphicx}
\usepackage{xcolor}
\usepackage{indentfirst}
\usepackage{amsfonts}
\usepackage{amsmath}
\usepackage{amsthm}
\usepackage{color}
\usepackage[backref=page]{hyperref}
\usepackage[capitalize]{cleveref}
\usepackage{verbatim}
\usepackage[makeroom]{cancel}
\usepackage{rotating}
\usepackage{float}
\usepackage{fullpage}
\usepackage{amssymb}
\usepackage{amsthm}
\usepackage{stmaryrd}
\usepackage{mathrsfs}
\usepackage{epsfig}
\usepackage{indentfirst}
\usepackage{framed}
\usepackage[numbers]{natbib}
\usepackage{color}
\usepackage{graphicx}
\usepackage{float}
\usepackage{bbm}
\usepackage{mathtools}
\usepackage{enumitem}
\usepackage{thmtools}
\usepackage{thm-restate}
\usepackage{longtable}
\usepackage{booktabs}
% \newtheorem{theorem}{Theorem}
\newtheorem{theorem}{Theorem}[section]
\newtheorem{question}[theorem]{Question}
\newtheorem{problem}[theorem]{Problem}
\newtheorem{lemma}[theorem]{Lemma}
\newtheorem{corollary}[theorem]{Corollary}
\newtheorem{definition}[theorem]{Definition}
\newtheorem{proposition}[theorem]{Proposition}
\newtheorem{remark}[theorem]{Remark}
\newtheorem{note}[theorem]{Note}
\newtheorem{fact}[theorem]{Fact}
\newtheorem{claim}[theorem]{Claim}
\newtheorem{assumption}[theorem]{Assumption}
\newtheorem{hypothesis}[theorem]{Hypothesis}
\newtheorem{conjecture}[theorem]{Conjecture}

\crefname{hypothesis}{hypothesis}{hypotheses}

% \numberwithin{theorem}{section}

\newtheorem*{claim*}{Claim}

% Change the parameter below to 0 to remove colors.
\def\colorful{1}

\ifnum\colorful=1
\newcommand{\violet}[1]{{\color{violet}{#1}}}
\newcommand{\orange}[1]{{\color{orange}{#1}}}
\newcommand{\blue}[1]{{\color{blue} {#1}}}
\newcommand{\red}[1]{{\color{red} {#1}}}
\newcommand{\green}[1]{{\color{green} {#1}}}
\newcommand{\pink}[1]{{\color{pink}{#1}}}
\newcommand{\gray}[1]{{\color{gray}{#1}}}

\fi
\ifnum\colorful=0
\newcommand{\violet}[1]{{{#1}}}
\newcommand{\orange}[1]{{{#1}}}
\newcommand{\blue}[1]{{{#1}}}
\newcommand{\red}[1]{{{#1}}}
\newcommand{\green}[1]{{{#1}}}
\fi


\newcommand{\bvee}{\bar{\vee}}
\newcommand{\bwedge}{\bar{\wedge}}
\newcommand{\lsem}{[\![}
\newcommand{\rsem}{]\!]}
% \newcommand{\eqdef}{\stackrel{\rm def}{=}}
\newcommand{\eqdef}{:=}
\newcommand{\defx}[1]{\emph{#1}}

\definecolor{brickred}{rgb}{0.8, 0.25, 0.33}
\definecolor{electriccrimson}{rgb}{1.0, 0.0, 0.25}
\definecolor{harvardcrimson}{rgb}{0.79, 0.0, 0.09}

\definecolor{camouflagegreen}{rgb}{0.47, 0.53, 0.42}
\definecolor{armygreen}{rgb}{0.29, 0.33, 0.13}
\definecolor{darkspringgreen}{rgb}{0.09, 0.45, 0.27}
\definecolor{forestgreen}{rgb}{0.13, 0.55, 0.13}

\definecolor{purpleheart}{rgb}{0.41, 0.21, 0.61}
\newcommand{\inote}[1]{\noindent  \textcolor{blue}{(\textbf{Igor:} #1\noindent)}}
% \newcommand{\bnote}[1]{\noindent \textcolor{purpleheart}{(\textbf{Bruno:} #1\noindent)}}
%\newcommand{\bnote}[1]{}
\newcommand{\bnote}[1]{\noindent  \textcolor{red}{(\textbf{Bruno:} #1\noindent)}}

\newcommand{\XORf}{\mathsf{XOR}^{\mathsf{fnc}}}
\newcommand{\XORg}{\mathsf{XOR}^{\mathsf{graph}}}
\newcommand{\AND}{\mathsf{AND}}
\newcommand{\IPn}{\mathsf{IP}_n}
\newcommand{\binary}{\mathsf{binary}}

% MATH UTILS

\newcommand{\eps}{\varepsilon}
\DeclarePairedDelimiter\ceil{\lceil}{\rceil}
\DeclarePairedDelimiter\floor{\lfloor}{\rfloor}
\newcommand{\one}{\mathbbm{1}}
\DeclareMathOperator{\supp}{supp}
% \DeclareMathOperator{\poly}{poly}
\newcommand{\poly}{\mathsf{poly}}
\DeclareMathOperator{\polylog}{polylog}
\newcommand\inner[2]{\left\langle #1, #2 \right\rangle}
\newcommand{\ovl}{\overline}


% COMPLEXITY CLASSES

% \def\P{{\calp}}
% \def\NP{{\caln\calp}}
% \def\P{{\text{P}}}
% \def\NP{{\text{NP}}}
\newcommand{\complexityfont}[1]{\mathsf{#1}}
\newcommand{\monclass}[1]{\mathsf{m}#1}
% \DeclareMathOperator{\CC}{\mathsf{CC}}
\def\mCC          {\monclass{\CC}}
\def\CC           {\complexityfont{CC}}
\def\TC           {\complexityfont{TC}}
\def\RNC          {\complexityfont{RNC}}
\def\mNC          {\monclass{\NC}}
\def\mTC          {\monclass{\TC}}
\def\EXPSPACE     {\complexityfont{EXPSPACE}}
\def\EXPTIME      {\complexityfont{EXPTIME}
\def\PSPACE       {\complexityfont{PSPACE}}}
\newcommand\NL           {\complexityfont{NL}}
\def\mNL          {\monclass{\NL}}
\def\BPP          {\complexityfont{BPP}}
\def\ZPP          {\complexityfont{ZPP}}
\def\PP           {\complexityfont{PP}}
\def\RP           {\complexityfont{RP}}
\def\coRP         {\complexityfont{coRP}}
\def\P            {\complexityfont{P}}
\def\Ppoly        {\complexityfont{P/poly}}
\def\mPpoly       {\monclass{\Ppoly}}
\def\NP           {\complexityfont{NP}}
\def\coNP         {\complexityfont{coNP}}
\def\EXP          {\complexityfont{EXP}}
\def\AC           {\complexityfont{AC}}
% \def\pL           {\oplus\text{-}\L}
% \def\pL           {\oplus\L}
% \def\pL           {\oplus\complexityfont{L}}
% \def\pL           {\oplus\complexityfont{LOGSPACE}}
\def\pL           {\oplus\complexityfont{L}}
\def\pLpoly       {\pL\complexityfont{/poly}}
\def\NLpoly       {\NL\complexityfont{/poly}}
\def\ACC          {\complexityfont{ACC}}
\def\MAJ          {\complexityfont{MAJ}}
\def\mAC          {\complexityfont{mAC}}
\def\mP           {\monclass{\P}}
\def\NC           {\complexityfont{NC}}
\def\NOT          {\complexityfont{NOT}}
\def\OR           {\complexityfont{OR}}
\def\NAND         {\complexityfont{NAND}}
\def\AND          {\complexityfont{AND}}
\def\SIZE         {\complexityfont{SIZE}}
\def\DEPTH        {\complexityfont{DEPTH}}
\def\mSIZE        {\monclass{\SIZE}}
\def\mDEPTH       {\monclass{\DEPTH}}
\renewcommand{\L} {\complexityfont{L}}
% \renewcommand{\L} {\complexityfont{DLOGSPACE}}
\newcommand{\KW}{\mathop{\text{KW}}}
\newcommand{\mKW}{\mathop{\text{KW}^+}}
\newcommand{\Lclass} {{\normalfont{\textbf{L}}}}
\newcommand{\SPF}[1]{\complexityfont{Span}_{\bbf_{#1}}}
\newcommand{\mSPF}[1]{\monclass{\complexityfont{Span}_{\bbf_{#1}}}}
\newcommand{\semMon}[1]{#1 \cap \monfns}
% \newcommand{\monfns}{\mathsf{mono}}
\newcommand{\monfns}{\mathsf{Mono}}
\newcommand{\GraphPpts}{\mathsf{GraphProperties}}
\newcommand{\HomPres}{\mathsf{HomPreserving}}
\newcommand{\txorsat}{\mathsf{3\mbox{-}XOR\mbox{-}SAT}}
\newcommand{\thornsat}{\mathsf{3\mbox{-}Horn\mbox{-}SAT}}
\newcommand{\tahornsat}{\mathsf{3\mbox{-}AntiHorn\mbox{-}SAT}}
% \newcommand{\txorsat}{\mathsf{3\mbox{-}\!XOR\mbox{-}\!SAT}}
\newcommand{\Gclean}[1]{{#1}_{\mathsf{clean}}}
\newcommand{\spoly}{\mathsf{/poly}}
\newcommand{\hornt}{\calh^3}
\newcommand{\ahornt}{\cala^3}
\newcommand{\NTISP}{\complexityfont{NTiSp}}



% CIRCUIT UTILS

\newcommand{\mDNF}{\monclass{\DNF}}
\newcommand{\mCNF}{\monclass{\CNF}}
\newcommand{\DNF}{\mathsf{DNF}}
\newcommand{\CNF}{\mathsf{CNF}}
\newcommand{\Cmonpref}{\mathsf{C}^{\mathsf{mon}}}
% \newcommand{\Cmon}[1]{\Cmonpref(#1)}
%\newcommand{\Cmon}[1]{\mathsf{mon{\text{-}}size}({#1})}
\newcommand{\Cmon}[1]{\mathsf{mSIZE}({#1})}
\newcommand{\Cman}[1]{\Cmonpref_{\land}(#1)}
\newcommand{\Cskw}[1]{\mathsf{C}_{\mathsf{skew}}(#1)}
\newcommand{\Ccsk}[1]{\mathsf{C}_{\mathsf{skew}}^\circ(#1)}
\newcommand{\Cska}[1]{\mathsf{C}^{\mathsf{skew}}_{\land}(#1)}
\newcommand{\Cszpref}{\mathsf{C}}
% \newcommand{\Cszpref}{\mathsf{C}_\mathsf{size}}
\newcommand{\Csz}[1]{\Cszpref(#1)}
\newcommand{\Csza}[1]{\Cszpref_{\land}(#1)}
% \newcommand{\Csza}[1]{\Cszpref^{\land}(#1)}
\newcommand{\apf}[1]{{#1}^{\cala}}
\newcommand{\GenFml}[1]{\calg_{#1}}
\newcommand{\mdepth}[1]{\mathsf{mDEPTH}({#1})}
\newcommand{\msize}[1]{\Cmon{#1}}

% DISCRETE COMPLEXITY

\newcommand{\bst}{\Gamma}

\newcommand{\gphgn}{\mathcal{G}_N}
\newcommand{\ghclq}{\mathcal{K}_{N,M}}
\newcommand{\boolg}{\mathcal{B}^{\mathsf{mon}}_{n}}
\newcommand{\boolgt}{\mathcal{B}_{2n}}
\newcommand{\mongen}[1]{\mathcal{B}^{\mathsf{mon}}_{#1}}
\newcommand{\blgen}[1]{\mathcal{B}_{#1}}
\newcommand{\Dand}[2]{D_{\cap}({#1} \mid {#2})}
\newcommand{\Dc}[2]{D(#1 \mid #2)}
\newcommand{\Dcg}[1]{\Dc{#1}{\gphgn}}
\newcommand{\Dcb}[1]{\Dc{#1}{\boolgt}}
\newcommand{\Dacg}[1]{\Dand{#1}{\gphgn}}
\newcommand{\Dacb}[1]{\Dand{#1}{\boolgt}}
\newcommand{\row}[1]{R^{#1}}
\newcommand{\col}[1]{C^{#1}}
\newcommand{\rwgc}[1]{\row{#1}_{\ovbr{G}}}
\newcommand{\clgc}[1]{\col{#1}_{\ovbr{G}}}
\newcommand{\flt}[1]{\mathfrak{F}_{\mathsf{mon}}(#1)}
\newcommand{\prj}[1]{[\lvert #1 \rvert]}
\newcommand{\rhobp}{\rho_{\text{BP}}}
\newcommand{\ftdef}[1]{\mathfrak{F}_{\mathsf{def}}(#1)}
\newcommand{\RS}{\mathsf{RS}}

%% COVER GRAPH

% \newcommand{\Vpairs}{\mathsf{Pairs}_U}
% \newcommand{\Vfilters}{\mathsf{Filters}_{A, U, \calb}}
% \newcommand{\Phedges}{\mathcal{E}_{A, U, \calb}}
\newcommand{\Vpairs}{V_\mathsf{pairs}}
\newcommand{\Vfilters}{V_\mathsf{filters}}
\newcommand{\Phedges}{\mathcal{E}}

%%% ABBREVIATIONS
% \renewcommand{\mb}{\bm}
\newcommand{\tu}{\textup}
% \newcommand{\ds}{\displaystyle}
%\newcommand{\ts}{\textstyle}
%\newcommand{\ul}{\underline}
\newcommand{\mr}{\mathrm}
\newcommand{\mi}{\mathit}
\newcommand{\mf}{\mathfrak}
\newcommand{\mc}{\mathcal}
\newcommand{\mb}{\mathbf}
\newcommand{\mbb}{\mathbb}
\newcommand{\mds}{\mathds}
% \newcommand{\scr}{\mathscr}
%\newenvironment{pf}{\begin{proof}}{\end{proof}}

% MATH OPERATORS

\def\sseq{\subseteq}
\newcommand{\set}[1]{\left\{#1\right\}}
\newcommand{\overbar}[1]{\mkern 1.5mu\overline{\mkern-1.5mu#1\mkern-1.5mu}\mkern 1.5mu}
\newcommand{\ovbr}{\overbar}
\def\sm{\setminus}
\def\var{\mathop{\text{\rm var}}\nolimits}
\newcommand{\ds}{\oplus}
\newcommand{\card}[1]{\lvert#1\rvert}
\newcommand{\abs}[1]{\left\lvert#1\right\rvert}
\newcommand{\hwt}[1]{\card{#1}_1}
% \newcommand{\cmpl}[1]{\bar{#1}}
% \newcommand{\cmpl}[1]{\ovbr{#1}}
% \newcommand{\cmpl}[1]{{#1}^{\mathsf{c}}}
\newcommand{\cmpl}[1]{#1^{\mathsf{c}}}
\newcommand{\brkt}[1]{\left\langle #1 \right\rangle}

% MATH OPERATORS (MATRIX)

% \newcommand{\tT}{\intercal}
% \newcommand{\tT}{\top}
% \newcommand{\tT}{T}
\newcommand{\tT}{\mathsf{T}}
\newcommand{\tr}[1]{#1^\tT}
\newcommand{\Diagld}[1]{\Diag \left( \ldm{A} \right)}

% BOOLEAN

%\newcommand{\CNF}{\mathsf{CNF}}
\newcommand{\blt}{\set{0,1}}
\newcommand{\vecp}{\vec{p}}
\newcommand{\vecq}{\vec{q}}
\newcommand{\vecr}{\vec{r}}
% \newcommand{\power}[1]{2^{#1}}
\newcommand{\power}[1]{\calp({#1})}
\DeclareMathOperator{\DT}{DT}
% \newcommand{\dtdepth}[1]{\DT_\text{\footnotesize depth}(#1)}
\newcommand{\dtdepth}[1]{\mathsf{DT}_\mathsf{depth}(#1)}
% \newcommand{\dtsize}[1]{\DT_\text{\footnotesize\rm size}(#1)}
\newcommand{\dtsize}[1]{\mathsf{DT}_\mathsf{size}(#1)}

% BOOLEAN FUNCTIONS

% \DeclareMathOperator{\clique}{Clique}
\newcommand{\clique}{\mathsf{Clique}}
\DeclareMathOperator{\cycle}{CYCLE}
\DeclareMathOperator{\thr}{THR}
% \newcommand{\vcprm}{k\mbox{-}\!\mathsf{Vertex\mbox{-}Cover}}
\newcommand{\kvc}{k\mbox{-}\!\mathsf{Vertex\mbox{-}Cover}}
\newcommand{\clqprm}[1]{#1\mbox{-}\!\clique}
\newcommand{\kclq}{k\mbox{-}\!\clique}
% \newcommand{\kclq}{\mathsf{CLIQUE}(n,k)}
% \newcommand{\kclq}{\mathsf{Clique}(n,k)}
\newcommand{\kmtc}{k\mbox{-}\!\Matching}
\newcommand{\Matching}{\mathsf{Matching}}
\newcommand{\BPM}{\mathsf{Bipartite}\text{-}\mathsf{Matching}}
\newcommand{\Parity}{\mathsf{Parity}}
\newcommand{\maj}{\mathsf{Majority}}
\newcommand{\MOD}{\mathsf{MOD}}
\newcommand{\andp}{\mathsf{POLY(q,s)}}
% \newcommand{\kclq}{\mathsf{Clique}_{k,n}}
\newcommand{\kcycl}{k$-$\cycle}
\newcommand{\MCSP}{\mathsf{MCSP}}
\newcommand{\stconn}{\mathsf{ST\text{-}CONN}}
\newcommand{\ustconn}{\mathsf{U\text{-}ST\text{-}CONN}}
\newcommand{\conn}{\mathsf{Connectivity}}
\newcommand{\GEN}{\mathsf{GEN}}
\newcommand{\SAT}{\mathsf{SAT}}
\newcommand{\dSAT}{2\text{-}\mathsf{SAT}}
\newcommand{\MtxSs}[1]{\llbracket#1\rrbracket}
\newcommand{\MSP}{\mathsf{MSP}}
\newcommand{\xor}{\oplus}
\newcommand{\bigxor}{\bigoplus}
\newcommand{\OddFactor}{\mathsf{OddFactor}}
\newcommand{\EqOne}{\mathsf{EqualONE}}
\newcommand{\EvenCC}{\mathsf{EvenCC}}
\newcommand{\BipEvenCC}{\mathsf{Bipartite}\text{-}\mathsf{EvenCC}}
\newcommand{\BipOddFactor}{\mathsf{Bipartite}\text{-}\OddFactor}


% CSP

\newcommand{\CSPSAT}{\mathsf{CSP}\text{-}\mathsf{SAT}}
% \newcommand{\COQf}[1]{\mathsf{COQ}\text{-}\mathsf{SAT}_{#1}}
\newcommand{\COQf}[1]{\mathsf{COQ}_{#1}}
\newcommand{\CSPf}[1]{\mathsf{CSP}\text{-}\mathsf{SAT}_{#1}}
% \newcommand{\CSPfn}[2]{\mathsf{CSP}\text{-}\mathsf{SAT}_{#1}^{#2}}
\newcommand{\CSPfn}[2]{\CSPf{#1}^{#2}}
\newcommand{\COQfn}[3]{\mathsf{COQ}_{#1}^{(#2,#3)}}
\newcommand{\CSP}{\mathsf{CSP}}
\newcommand{\COQ}{\mathsf{COQ}}
\newcommand{\Pol}{\mathsf{Pol}}
\newcommand{\Inv}{\mathsf{Inv}}
\newcommand{\ccln}[1]{\brkt{#1}}
\newcommand{\Cstr}[1]{\mathsf{Constraints}(#1)}
\newcommand{\Cstrn}[2]{\mathsf{Constraints}_{#1}(#2)}
\newcommand{\proj}{\mathsf{proj}}
\newcommand{\dual}{\mathsf{dual}}
\newcommand{\acmred}{\leq^{\AC^0}_m\!}
\newcommand{\acmcmpl}{$\acmred$\!-complete}
\newcommand{\mnlmred}{\leq^{\mNL}_m}
\newcommand{\mprojred}{\leq^{\mathsf{mProj}}_m}
\newcommand{\mormred}{\leq^{\mathsf{mOR}}_m}
\newcommand{\clonefont}[1]{\mathrm{#1}}


% COMMUNICATION COMPLEXITY

\newcommand{\PLS}{\mathsf{PLS}}
\newcommand{\Root}{\emptyset}

% LIFTING

% \newcommand{\SrchF}[1]{\mathsf{Search}_{#1}}
\newcommand{\SrchF}[1]{\mathsf{Search}(#1)}
\newcommand{\Clauses}[1]{\mathsf{Clauses}(#1)}


% CIRCUIT CLASSES

% \newcommand{\AC}{\mathsf{AC}}
% \newcommand{\NC}{\mathsf{NC}}

% FUSION METHOD

\DeclareMathOperator{\Gen}{\mathsf{GEN}}

% PROBABILITY STUFF

% \newcommand{\probabilityfont}[1]{\mathsf{#1}}
\newcommand{\Bin}{\text{Bin}}
\newcommand{\dstr}{\sim}
\newcommand{\probabilityfont}[1]{\mb{#1}}
\newcommand{\rnd}[1]{\probabilityfont{#1}}
\def\Pr{\mathop{{}\mathbb{P}}}
% \def\Pr{\mathop{\text{\rm Pr}}}
% \def\Exp{\mathop{\bbe}limits}
% \def\Exp{\mathop{\bbe}}
\def\Exp{\mathop{{}\mathbb{E}}}
\def\var{\mathop{\text{\rm var}}\nolimits}
\def\pcritk{p_{\rm c}^{K_k}}
\def\pcritl{p_{\rm c}^{C_k}}
\newcommand{\inuf}{\in_U}
\newcommand{\cktsz}{\calc}
% \newcommand{\inuf}{\in_calu}
\newcommand{\rndY}{\probabilityfont{Y}}
\newcommand{\rndN}{\probabilityfont{N}}
\newcommand{\rndW}{\probabilityfont{W}}
\newcommand{\rndP}{\probabilityfont{P}}
\newcommand{\rndS}{\probabilityfont{S}}
\newcommand{\rndT}{\probabilityfont{T}}
\newcommand{\rndX}{\probabilityfont{X}}
\newcommand{\rndG}{\probabilityfont{G}}
\newcommand{\rndH}{\probabilityfont{H}}
\newcommand{\rndF}{\probabilityfont{F}}
\newcommand{\rndK}{\probabilityfont{K}}
\newcommand{\rndc}{\probabilityfont{c}}
\newcommand{\rndx}{\probabilityfont{x}}
\newcommand{\flws}{\sim}


\def\bba{{\mathbb A}}
\def\bbb{{\mathbb B}}
\def\bbc{{\mathbb C}}
\def\bbd{{\mathbb D}}
\def\bbe{{\mathbb E}}
\def\bbf{{\mathbb F}}
\def\bbg{{\mathbb G}}
\def\bbh{{\mathbb H}}
\def\bbi{{\mathbb I}}
\def\bbj{{\mathbb J}}
\def\bbk{{\mathbb K}}
\def\bbl{{\mathbb L}}
\def\bbm{{\mathbb M}}
\def\bbn{{\mathbb N}}
\def\bbo{{\mathbb O}}
\def\bbp{{\mathbb P}}
\def\bbq{{\mathbb Q}}
\def\bbr{{\mathbb R}}
\def\bbs{{\mathbb S}}
\def\bbt{{\mathbb T}}
\def\bbu{{\mathbb U}}
\def\bbv{{\mathbb V}}
\def\bbw{{\mathbb W}}
\def\bbx{{\mathbb X}}
\def\bby{{\mathbb Y}}
\def\bbz{{\mathbb Z}}

\def\cala{{\mathcal A}}
\def\calb{{\mathcal B}}
\def\calc{{\mathcal C}}
\def\cald{{\mathcal D}}
\def\cale{{\mathcal E}}
\def\calf{{\mathcal F}}
\def\calg{{\mathcal G}}
\def\calh{{\mathcal H}}
\def\cali{{\mathcal I}}
\def\calj{{\mathcal J}}
\def\calk{{\mathcal K}}
\def\call{{\mathcal L}}
\def\calm{{\mathcal M}}
\def\caln{{\mathcal N}}
\def\calo{{\mathcal O}}
\def\calp{{\mathcal P}}
\def\calq{{\mathcal Q}}
\def\calr{{\mathcal R}}
\def\cals{{\mathcal S}}
\def\calt{{\mathcal T}}
\def\calu{{\mathcal U}}
\def\calv{{\mathcal V}}
\def\calw{{\mathcal W}}
\def\calx{{\mathcal X}}
\def\caly{{\mathcal Y}}
\def\calz{{\mathcal Z}}

\def\scra{{\mathscr A}}
\def\scrb{{\mathscr B}}
\def\scrc{{\mathscr C}}
\def\scrd{{\mathscr D}}
\def\scre{{\mathscr E}}
\def\scrf{{\mathscr F}}
\def\scrg{{\mathscr G}}
\def\scrh{{\mathscr H}}
\def\scri{{\mathscr I}}
\def\scrj{{\mathscr J}}
\def\scrk{{\mathscr K}}
\def\scrl{{\mathscr L}}
\def\scrm{{\mathscr M}}
\def\scrn{{\mathscr N}}
\def\scro{{\mathscr O}}
\def\scrp{{\mathscr P}}
\def\scrq{{\mathscr Q}}
\def\scrr{{\mathscr R}}
\def\scrs{{\mathscr S}}
\def\scrt{{\mathscr T}}
\def\scru{{\mathscr U}}
\def\scrv{{\mathscr V}}
\def\scrw{{\mathscr W}}
\def\scrx{{\mathscr X}}
\def\scry{{\mathscr Y}}
\def\scrz{{\mathscr Z}}
\newcommand{\headsubsec}[1]{\paragraph{#1.}}
\newcommand{\headsubsecline}[1]{\headsubsec{#1}}
\newcommand{\gaphyp}{monotone gap conjecture}
\newcommand{\gaphyps}{monotone gap conjectures}
\newcommand{\ckthyp}{monotone circuit gap conjecture}
\newcommand{\fmlhyp}{monotone formula gap conjecture}
\usepackage[labelfont=bf]{caption}
\makeatletter
%% postlattice.sty
%% Copyright 2016 by Gernot Salzer, salzer@logic.at
%
% This work may be distributed and/or modified under the
% conditions of the LaTeX Project Public License, either version 1.3c
% of this license or (at your option) any later version.
% The latest version of this license is in
%   http://www.latex-project.org/lppl.txt
% and version 1.3 or later is part of all distributions of LaTeX
% version 2005/12/01 or later.
%
% This work has the LPPL maintenance status `maintained',
% the current maintainer being Gernot Salzer.

% postlattice.sty ... package for drawing and decorating Post's lattice
%   of Boolean functional/relational clones

\NeedsTeXFormat{LaTeX2e}
\ProvidesPackage{postlattice}[2016/08/18 Post's lattice of Boolean clones]
\RequirePackage{tikz}
\usetikzlibrary{patterns}

\newcommand\postset{\pgfqkeys{/post}}

\postset
 {unknownkeys/.style={},
  .unknown/.code=%
   {\let\@n\pgfkeyscurrentname
    \let\@v\pgfkeyscurrentvalue
    \ifx#1\pgfkeysnovalue
      \postset{unknownkeys/.append style/.expand once={\@n}}%
    \else
      \postset{unknownkeys/.append style/.expand twice={\expandafter\@n\expandafter=\@v}}%
    \fi
   },
  labels/.is choice,
  labels/.value required,
  labels/none/.code={\postlabelsnone},
  labels/internal/.code={\postlabelsinternal},
  labels/wagner/.code={\postlabelswagner},
  labels/iwagner/.code={\postlabelsiwagner},
  labels/post/.code={\postlabelspost},
  labels/poekal/.code={\postlabelspoekal},
  labels/wikipedia/.code={\postlabelswikipedia},
  layout/.is choice,
  layout/.value required,
  layout/circles/.code={\postlayoutcircles},
  layout/circlesX/.code={\postlayoutcirclesX},
  layout/bullets/.code={\postlayoutbullets},
  layout/bulletsX/.code={\postlayoutbulletsX},
  layout/lau/.code={\postlayoutlau},
  /handlers/.lstyle/.code=
    \edef\pgfkeyscurrentkey{\pgfkeyscurrentpath}%
    \pgfkeys@split@path
    \pgfkeysalso{lstyle/\pgfkeyscurrentname/.style={#1}},
  lstyle/.unknown/.code=
    \expandafter\droplast\expandafter{\pgfkeyscurrentkey}%
    \pgfkeysalso{\the\lastdropped},
  /handlers/.cstyle/.code=
    \edef\pgfkeyscurrentkey{\pgfkeyscurrentpath}%
    \pgfkeys@split@path
    \pgfkeysalso{cstyle/\pgfkeyscurrentname/.style={#1}},
  cstyle/.unknown/.code=
    \expandafter\droplast\expandafter{\pgfkeyscurrentkey}%
    \pgfkeysalso{\the\lastdropped},
  /handlers/.pos/.code=
    \edef\pgfkeyscurrentkey{\pgfkeyscurrentpath}%
    \pgfkeys@split@path
    \pgfkeysalso{pos/\pgfkeyscurrentname={#1}},
  labelformat/.code=
    \pgfkeysalso{formatlabel/.code={#1}},
  labelformat={{\scriptsize$\mathrm{#1}$}},
  cshape/circle/.style={circle,draw,inner sep=0,outer sep=0,minimum width=6mm,fill=white},
  cshape/bullet/.style={circle,draw,inner sep=0,outer sep=0,minimum width=1mm,fill=black}
 }
\pgfkeysdef{/post/lstyle/}{}
\pgfkeysdef{/post/cstyle/}{}

\newtoks\lastdropped
\newcommand\droplast[1]{\lastdropped{}\@droplast#1\@end}
\def\@droplast#1#2%
 {\ifx\@end#2%
    \let\@@droplast\@@@droplast
  \else
    \expandafter\lastdropped\expandafter{\the\lastdropped#1}%
    \let\@@droplast\@droplast
  \fi
  \@@droplast#2%
 }
\def\@@@droplast#1{}

\newcommand\postnodes
 {\pn{C0}\pn{C1}\pn{C2}\pn{Cx}\pn{D1}\pn{D2}\pn{Dx}%
  \pn{E0}\pn{E1}\pn{E2}\pn{Ex}\pn{I0}\pn{I1}\pn{I2}\pn{Ix}%
  \pn{L0}\pn{L1}\pn{L2}\pn{L3}\pn{Lx}\pn{M0}\pn{M1}\pn{M2}\pn{Mx}%
  \pn{N2}\pn{Nx}\pn{R0}\pn{R1}\pn{R2}\pn{Rx}%
  \pn{S00}\pn{S01}\pn{S02}\pn{S0x}\pn{S10}\pn{S11}\pn{S12}\pn{S1x}%
  \pn{S200}\pn{S201}\pn{S202}\pn{S20x}\pn{S210}\pn{S211}\pn{S212}\pn{S21x}%
  \pn{S300}\pn{S301}\pn{S302}\pn{S30x}\pn{S310}\pn{S311}\pn{S312}\pn{S31x}%
  \pn{V0}\pn{V1}\pn{V2}\pn{Vx}%
 }

\newcommand\postlabelsnone
 {\def\pn##1{\pgfkeyssetvalue{/post/##1}{}}%
  \postnodes
 }

\newcommand\postlabelsinternal
 {\def\pn##1{\pgfkeyssetvalue{/post/##1}{##1}}%
  \postnodes
 }

\newcommand\postlabelswikipedia
 {\postset
   {C0=,C1=,C2=,Cx=,
    D1=DP,D2=DM,Dx=D,
    E0=\Lambda P_0,E1=\Lambda P_1,E2=\Lambda P,Ex=\Lambda,
    I0=UP_0,I1=UP_1,I2=\bot,Ix=UM,
    L0=AP_0,L1=AP_1,L2=AP,L3=AD,Lx=A,
    M0=MP_0,M1=MP_1,M2=MP,Mx=M,
    N2=UD,Nx=U,
    R0=P_0,R1=P_1,R2=P,Rx=\top,
    S00=MPT_1^\infty,S01=MT_1^\infty,S02=PT_1^\infty,S0x=T_1^\infty,
    S10=MPT_0^\infty,S11=MT_0^\infty,S12=PT_0^\infty,S1x=T_0^\infty,
    S200=MPT_1^2,S201=MT_1^2,S202=PT_1^2,S20x=T_1^2,
    S210=MPT_0^2,S211=MT_0^2,S212=PT_0^2,S21x=T_0^2,
    S300=MPT_1^3,S301=MT_1^3,S302=PT_1^3,S30x=T_1^3,
    S310=MPT_0^3,S311=MT_0^3,S312=PT_0^3,S31x=T_0^3,
    V0=VP_0,V1=VP_1,V2=VP,Vx=V
   }%
 }

\newcommand\postlabelspost
 {\postset
   {C0=O_3,C1=O_2,C2=,Cx=O_7,
    D1=D_1,D2=D_2,Dx=D_3,
    E0=P_3,E1=P_5,E2=P_1,Ex=P_6,
    I0=O_6,I1=O_5,I2=O_1,Ix=O_8,
    L0=L_3,L1=L_2,L2=L_4,L3=L_5,Lx=L_1,
    M0=A_3,M1=A_2,M2=A_4,Mx=A_1,
    N2=O_4,Nx=O_9,
    R0=C_3,R1=C_2,R2=C_4,Rx=C_1,
    S00=F_2^\infty,S01=F_3^\infty,S02=F_1^\infty,S0x=F_4^\infty,
    S10=F_6^\infty,S11=F_7^\infty,S12=F_5^\infty,S1x=F_8^\infty,
    S200=F_2^2,S201=F_3^2,S202=F_1^2,S20x=F_4^2,
    S210=F_6^2,S211=F_7^2,S212=F_5^2,S21x=F_8^2,
    S300=F_2^3,S301=F_3^3,S302=F_1^3,S30x=F_4^3,
    S310=F_6^3,S311=F_7^3,S312=F_5^3,S31x=F_8^3,
    V0=S_5,V1=S_3,V2=S_1,Vx=S_6
   }
 }

\newcommand\postlabelspoekal
 {\postlabelspost
  \postset{M0=M_3,M1=M_2,M2=M_4,Mx=M_1}%
 }

\newcommand\postlabelswagner
 {\def\pn##1%
   {\edef\l{\expandafter\@postlabelwagner##1{}{}{}{}\@end}%
    \pgfkeyslet{/post/##1}{\l}%
   }%
  \postnodes
  \postset{Rx=BF}%
 }

\newcommand\postlabelsiwagner
 {\def\pn##1%
   {\edef\l{i\expandafter\@postlabelwagner##1{}{}{}{}\@end}%
    \pgfkeyslet{/post/##1}{\l}%
   }%
  \postnodes
  \postset{I2=BR,Rx=iBF}
 }

\newcommand\@postlabelwagner{}
\def\@postlabelwagner#1#2#3#4#5\@end
 {\ifx S#1%
    \ifnum#2>1%
      S^{#2}_{#3\ifx x#4\else#4\fi}%
    \else
      S_{#2\ifx x#3\else#3\fi}%
    \fi
  \else
    #1\ifx x#2\else_{#2}\fi%
  \fi
 }

\newcommand\labelposcenter{}
\def\labelposcenter{c:c}
\newcommand\postpositionscircle
 {\def\pn##1{\pgfkeyslet{/post/pos/##1}{\labelposcenter}}
  \postnodes
 }

\newcommand\postpositionsbullet
 {\postset{pos/.cd,
    C0=0:0.1, C1=180:0.1, C2=-90:0.1, Cx=-90:0.2,
    D1=-10:0.1, D2=-10:0.1, Dx=10:0.1,
    E0=0:0.1, E1=180:0.1, E2=-65:0.15, Ex=55:0.2,
    I0=-37:0.1, I1=-143:0.1, I2=-90:0.2, Ix=-103:0.25,
    L0=0:0.13, L1=180:0.13, L2=-63:0.15, L3=0:0.15, Lx=75:0.15,
    M0=18:0.1, M1=-108:0.25, M2=-90:0.2, Mx=120:0.2,
    N2=170:0.1, Nx=-165:0.1,
    R0=90:0.2, R1=90:0.2, R2=90:0.2, Rx=90:0.2,
    S00=0:0.15, S01=18:0.1, S02=-162:0.15, S0x=180:0.15,
    S10=180:0.15, S11=162:0.1, S12=-18:0.15, S1x=0:0.15,
    S200=90:0.2, S201=90:0.2, S202=-162:0.1, S20x=180:0.15,
    S210=90:0.2, S211=90:0.2, S212=-18:0.1, S21x=0:0.15,
    S300=0:0.15, S301=18:0.1, S302=-162:0.1, S30x=180:0.15,
    S310=180:0.15, S311=162:0.1, S312=-18:0.1, S31x=0:0.15,
    V0=0:0.1, V1=180:0.1, V2=-125:0.15, Vx=125:0.2
   }
 }

\newcommand\postpositionslau
 {\postset{pos/.cd,
    C0=-90:0.1, C1=-90:0.1, C2=-90:0.1, Cx=-90:0.1,
    D1=0:0.1, D2=-45:0.1, Dx=0:0.08,
    E0=-90:0.1, E1=-90:0.1, E2=-90:0.1, Ex=-90:0.1,
    I0=-34:0.1, I1=-156:0.1, I2=-35:0.1, Ix=130:0.13,
    L0=-50:0.15, L1=-130:0.15, L2=-10:0.3, L3=180:0.3, Lx=135:0.15,
    M0=35:0.12, M1=165:0.2, M2=-90:0.1, Mx=35:0.12,
    N2=-125:0.1, Nx=160:0.1,
    R0=90:0.13, R1=90:0.13, R2=-140:0.15, Rx=90:0.1,
    S00=0:0.1, S01=18:0.1, S02=-162:0.1, S0x=180:0.1,
    S10=180:0.1, S11=162:0.1, S12=-18:0.1, S1x=0:0.1,
    S200=70:0.15, S201=90:0.1, S202=-162:0.1, S20x=180:0.1,
    S210=110:0.15, S211=90:0.1, S212=-18:0.1, S21x=0:0.1,
    S300=0:0.1, S301=18:0.1, S302=-162:0.1, S30x=180:0.1,
    S310=180:0.1, S311=162:0.1, S312=-18:0.1, S31x=0:0.1,
    V0=-90:0.1, V1=-90:0.1, V2=-90:0.1, Vx=-90:0.1
  }
 }

\newcommand\postlayoutcircles
 {\postpositionscircle
  \postset{cshape/.style={/post/cshape/circle}}
  \let\postlayout\postlayoutA
 }

\newcommand\postlayoutcirclesX
 {\postpositionscircle
  \postset{cshape/.style={/post/cshape/circle}}
  \let\postlayout\postlayoutB
 }

\newcommand\postlayoutbullets
 {\postpositionsbullet
  \postset{cshape/.style={/post/cshape/bullet}}
  \let\postlayout\postlayoutA
 }

\newcommand\postlayoutbulletsX
 {\postpositionsbullet
  \postset{cshape/.style={/post/cshape/bullet}}
  \let\postlayout\postlayoutB
 }

\newcommand\postlayoutlau
 {\postpositionslau
  \postset{cshape/.style={/post/cshape/bullet}}
  \let\postlayout\postlayoutC
 }

\newif\if@projections
\newcommand\postlattice
 {\@ifstar
   {\@projectionstrue\@postlattice}%
   {\@projectionsfalse\@postlattice}%
 }
\newcommand\@postlattice[1][]%
 {{\postset{#1}%
  \begin{tikzpicture}[/post/unknownkeys,relative]
    \postlayout
  \end{tikzpicture}%
 }}

\newcommand\postlayoutA
 {\def\xE  { 3.35}% x-offset E- and V-clones
  \def\xS  { 5.45}% x-offset S-, S^2-, and S^2-clones
  \def\yNc { 2.25}% y-offset N2
  \def\yNx { 3.0}%  y-offset Nx
  \def\yL  { 4.5}%  y-offset L-clones
  \def\yDc { 6.35}% y-offset D2
  \def\yS  { 7.15}% y-offset S-clones
  \def\yDb { 7.15}% y-offset D1
  \def\yDx { 7.95}% y-offset Dx
  \def\ySd { 9.0}%  y-offset S^3-clones
  \def\ySc {10.15}% y-offset S^2-clones
  \def\yM  {12.4}%  y-offset M-clones
  \def\yR  {15.4}%  y-offset R-clones
  \def\ydp {0.75}%  diamond/parallelogram height
  \def\xdia{1.15}%  diamond small width
  \def\xDia{1.5}%   diamond width
  \def\xpar{0.575}% parallelogram small width
  \def\xPar{1.33}%  parallelogram width
  \@postlayoutAB
 }

\newcommand\postlayoutB
 {\def\xE  { 3.4}% x-offset E- and V-clones
  \def\xS  { 6.0}% x-offset S-, S^2-, and S^2-clones
  \def\yNc { 2.5}% y-offset N2
  \def\yNx { 3.3}% y-offset Nx
  \def\yL  { 5.0}% y-offset L-clones
  \def\yDc { 6.7}% y-offset D2
  \def\yS  { 7.0}% y-offset S-clones
  \def\yDb { 7.5}% y-offset D1
  \def\yDx { 8.3}% y-offset Dx
  \def\ySd { 8.5}% y-offset S^3-clones
  \def\ySc { 9.5}% y-offset S^2-clones
  \def\yM  {11.5}% y-offset M-clones
  \def\yR  {14.0}% y-offset R-clones
  \def\ydp { 0.8}% diamond/parallelogram height
  \def\xdia{ 1.2}% diamond small width
  \def\xDia{ 1.2}% diamond width 
  \def\xpar{ 0.6}% parallelogram small width
  \def\xPar{ 1.8}% parallelogram width
  \@postlayoutAB
 }

\newcommand\@postlayoutAB
 {\@postABa{   0}{  0}{I}
  \clone   {   0,\yNc}{N2}
  \clone   {   0,\yNx}{Nx}
  \@postABa{-\xE}{\yL}{V}
  \@postABa{   0}{\yL}{L}
  \clone   {   0,\yL }{L3}
  \@postABa{ \xE}{\yL }{E}
  \clone   {   0,\yDc}{D2}
  \@postABc{-\xS}{\yS }{S0}
  \@postABd{ \xS}{\yS }{S1}
  \clone   {   0,\yDb}{D1}
  \clone   {   0,\yDx}{Dx}
  \@postABc{-\xS}{\ySd}{S30}
  \@postABd{ \xS}{\ySd}{S31}
  \@postABc{-\xS}{\ySc}{S20}
  \@postABd{ \xS}{\ySc}{S21}
  \@postABb{   0}{\yM }{M}
  \@postABb{   0}{\yR }{R}
  \draw (I2) -- (V2) -- (S00)
        (I2) -- (E2) -- (S10)
        (I1) -- (E1) -- (M1) -- (R1)
        (I1) -- (V1) -- (S01)
        (I0) -- (V0) -- (M0) -- (R0)
        (I0) -- (E0) -- (S11)
        (Ix) -- (Vx) -- (Mx)
        (Ix) -- (Ex) -- (Mx)
        (N2) -- (Nx)
        (L2) -- (L3) -- (Lx)
        (D2) -- (D1) -- (Dx)
        (D2) -- (S200)
        (D2) -- (S210)
        (S300) -- (S200) -- (M2)
        (S301) -- (S201) -- (M1)
        (S302) -- (S202) -- (R2)
        (S30x) -- (S20x) -- (R1)
        (S310) -- (S210) -- (M2)
        (S311) -- (S211) -- (M0)
        (S312) -- (S212) -- (R2)
        (S31x) -- (S21x) -- (R0);
  \draw (I2) to[out=27,in=140] (L2);
  \draw (I2) to[out=20,in=150] (N2);
  \draw (I2) to[out=-15,in=-150] (D2);
  \draw (I1) to[out=20,in=160] (L1);
  \draw (I0) to[out=-20,in=-160] (L0);
  \draw (Ix) to[out=-20,in=-135] (Nx);
  \draw (N2) to[out=-40,in=-140] (L3);
  \draw (Nx) to[out=40,in=140] (Lx);
  \draw (L2) to[out=40,in=140] (D1);
  \draw (L1) to[out=35,in=150] (R1);
  \draw (L3) to[out=-40,in=-140] (Dx);
  \draw (L0) to[out=-35,in=-150] (R0);
  \draw (Lx) to[out=40,in=140] (Rx);
  \draw (D1) to[out=30,in=147] (R2);
  \draw (Dx) to[out=-10,in=-150] (Rx);
  \draw (M2) to[out=33,in=153] (R2);
  \draw (Mx) to[out=-40,in=-160] (Rx);
  \draw[dashed] (S00) -- (S300)
                (S01) -- (S301)
                (S02) -- (S302)
                (S0x) -- (S30x)
                (S10) -- (S310)
                (S11) -- (S311)
                (S12) -- (S312)
                (S1x) -- (S31x);
  \if@projections
    \@postABa{   0}{-\yNc}{C}
    \draw (C0)--(I0) (C1)--(I1);
    \draw (Cx) to[out=-40,in=-140] (Ix);
    \draw (C2) to[out=40,in=140] (I2);
  \fi
 }
\newcommand\@postABa[3]%
 {\clone{#1,#2+\ydp}{#3x}
  \clone{#1+\xdia,#2}{#30}
  \clone{#1-\xdia,#2}{#31}
  \clone{#1,#2-\ydp}{#32}
  \draw (#3x) -- (#30) -- (#32) -- (#31) -- (#3x);   
 }
\newcommand\@postABb[3]%
 {\clone{#1,#2+\ydp}{#3x}
  \clone{#1+\xDia,#2}{#30}
  \clone{#1-\xDia,#2}{#31}
  \clone{#1,#2-\ydp}{#32}
  \draw (#3x) -- (#30) -- (#32) -- (#31) -- (#3x);   
 }
\newcommand\@postABc[3]%
 {\clone{#1-\xPar,#2+\ydp}{#3x}
  \clone{#1+\xPar,#2-\ydp}{#30}
  \clone{#1+\xpar,#2}{#31}
  \clone{#1-\xpar,#2}{#32}
  \draw (#3x) -- (#31) -- (#30) -- (#32) -- (#3x);   
 }
\newcommand\@postABd[3]%
 {\clone{#1+\xPar,#2+\ydp}{#3x}
  \clone{#1-\xPar,#2-\ydp}{#30}
  \clone{#1-\xpar,#2}{#31}
  \clone{#1+\xpar,#2}{#32}
  \draw (#3x) -- (#31) -- (#30) -- (#32) -- (#3x);   
 }

\newcommand\postlayoutC
 {\@postCa{ 0.00}{0.00}{I}
  \@postCb{-3.50}{0.50}{V}
  \clone  {-0.17,0.50}{N2}
  \@postCa{ 3.00}{0.50}{E}
  \clone  {-0.50,1.50}{Nx}
  \@postCc{-3.50}{2.00}{S0}
  \@postCa{ 0.00}{2.00}{L}
  \@postCd{ 3.00}{2.00}{S1}
  \clone  {-0.30,2.50}{L3}
  \@postCc{-3.50}{4.00}{S30}
  \clone  { 0.50,4.00}{D2}
  \@postCd{ 3.00}{4.00}{S31}
  \clone  { 0.00,4.55}{D1}
  \clone  {-0.37,5.00}{Dx}
  \@postCc{-3.50}{5.50}{S20}
  \@postCd{ 3.00}{5.50}{S21}
  \@postCa{ 0.20}{6.50}{M}
  \@postCa{ 0.00}{7.00}{R}
  \draw (I2) -- (E2) -- (S10)
        (I2) -- (D2) -- (D1) -- (Dx)
        (I2) -- (L2) -- (D1) -- (R2)
        (I2) -- (N2) -- (L3) -- (Dx) -- (Rx)
        (I2) -- (V2) -- (S00)
        (I0) -- (E0) -- (S11) 
        (I0) -- (L0) -- (R0)
        (I0) -- (V0)
        (L2) -- (L3) -- (Lx)
        (M0) -- (R0)
        (Mx) -- (Rx)
        (I1) -- (V1) -- (S01)
        (I1) -- (L1) -- (R1)
        (I1) -- (E1)
        (Ix) -- (Nx) -- (Lx) -- (Rx)
        (Ix) -- (Ex)
        (Ix) -- (Vx)
        (N2) -- (Nx)
        (D2) -- (S200)
        (D2) -- (S210)
        (S300) -- (S200) -- (M2) -- (R2)
        (S301) -- (S201) -- (M1) -- (R1)
        (S302) -- (S202) -- (R2)
        (S30x) -- (S20x) -- (R1)
        (S310) -- (S210) -- (M2)
        (S311) -- (S211) -- (M0)
        (S312) -- (S212) -- (R2)
        (S31x) -- (S21x) -- (R0);
  \draw[dashed]
        (S00) -- (S300)
        (S01) -- (S301)
        (S02) -- (S302)
        (S0x) -- (S30x)
        (S10) -- (S310)
        (S11) -- (S311)
        (S12) -- (S312)
        (S1x) -- (S31x);
  \draw (V0) to[out=20,in=160] (M0);
  \draw (Vx) to[out=20,in=160] (Mx);
  \draw (E1) to[out=-20,in=-160] (M1);
  \draw (Ex) to[out=-20,in=-160] (Mx);
  \if@projections
    \@postCa{0}{-1.5}{C};
    \draw (Cx) -- (Ix)
          (C0) -- (I0)
          (C1) -- (I1)
          (C2) -- (I2);
  \fi
 }
\newcommand\@postCa[3]%
 {\clone{#1-0.5,#2+1.0}{#3x}
  \clone{#1+1.0,#2+0.5}{#30}
  \clone{#1-1.5,#2+0.5}{#31}
  \clone{#1,#2}{#32}
  \draw (#3x) -- (#30) -- (#32) -- (#31) -- (#3x);
 }
\newcommand\@postCb[3]%
 {\clone{#1+0.5,#2+1.0}{#3x}
  \clone{#1+1.5,#2+0.5}{#30}
  \clone{#1-1.0,#2+0.5}{#31}
  \clone{#1,#2}{#32}
  \draw (#3x) -- (#30) -- (#32) -- (#31) -- (#3x);
 }
\newcommand\@postCc[3]%
 {\clone{#1-3.0,#2+1.0}{#3x}
  \clone{#1-1.0,#2+0.7}{#31}
  \clone{#1-2.0,#2+0.3}{#32}
  \clone{#1,#2}{#30}
  \draw (#3x) -- (#31) -- (#30) -- (#32) -- (#3x);
 }
\newcommand\@postCd[3]%
 {\clone{#1+3.0,#2+1.0}{#3x}
  \clone{#1+1.0,#2+0.7}{#31}
  \clone{#1+2.0,#2+0.3}{#32}
  \clone{#1,#2}{#30}
  \draw (#3x) -- (#31) -- (#30) -- (#32) -- (#3x);
 }

\newcommand\clone[2]%
 {\edef\labelpos{\pgfkeysvalueof{/post/pos/#2}}%
  \draw[x=1cm,y=1cm] (#1) node[name=#2,/post/cshape,/post/cstyle/#2] {}
    \ifx\labelpos\labelposcenter
      \centeredlabel{#2}%
    \else
      \expandafter\anchoredlabel\expandafter(\labelpos){#2}%
    \fi
  ;%
 }
\newcommand\centeredlabel[1]%
 {node[inner sep=0,outer sep=0,/post/lstyle/#1]
  {\postset{formatlabel=\pgfkeysvalueof{/post/#1}}}
 }
\newcommand\anchoredlabel{}
\def\anchoredlabel(#1:#2)#3%
 {+(#1:#2)
  node[anchor at angle=#1+180,inner sep=0,outer sep=0,/post/lstyle/#3]
  {\postset{formatlabel=\pgfkeysvalueof{/post/#3}}}
 }
\tikzset
 {anchor at angle/.code=
   {\pgfpointtransformed{\pgfpointpolar{#1}{1pt}}
    \pgfmathanglebetweenpoints{\pgfqpoint{-\pgf@x}{-\pgf@y}}{\pgfpointorigin}
    \def\tikz@anchor{\pgfmathresult}
%    \show\pgfmathresult
   }
 }

\postset{layout=circles,labels=none}

\makeatother
\postset{scale=-1,labels=wikipedia}
\newcommand\postlegend[1]%
 {\raisebox{-0.45ex}{\tikz\node[circle,draw,minimum size=2ex,#1]{};}}
\tikzset{NA/.style = {fill = white}}
\tikzset{pL/.style = {fill = green!70}}
\tikzset{NL/.style = {fill = yellow!70}}
\tikzset{L/.style = {fill = orange!80}}
\tikzset{LN/.style = {fill = red!70}}
\tikzset{NP/.style = {fill = cyan!80}}
\tikzset{P/.style = {fill = magenta!70}}
\definecolor{mikadoyellow}{rgb}{1.0, 0.77, 0.05}
\definecolor{cyan(process)}{rgb}{0.0, 0.72, 0.92}
\definecolor{applegreen}{rgb}{0.55, 0.71, 0.0}

\crefname{originalreference}{original reference}{original references}
\Crefname{originalreference}{Original reference}{Original references}
\makeatletter
\newlabel{s:constant-depth}{{3}{source}{}{sourcepointer-1}{}}
\newlabel{s:constant-depth@cref}{{[originalreference][0][]3}{[1][0][]source}{}{}{}}
\makeatother
\crefname{originalreference}{original reference}{original references}
\Crefname{originalreference}{Original reference}{Original references}
\makeatletter
\newlabel{s:padding_graph}{{3.2}{source}{}{sourcepointer-2}{}}
\newlabel{s:padding_graph@cref}{{[originalreference][0][]3.2}{[1][0][]source}{}{}{}}
\makeatother
\title{Unconditional monotone size and depth dichotomies for Boolean-valued constraint satisfaction}
\author{Bruno P. Cavalar, Igor C. Oliveira}
\date{doi:10.4230/LIPIcs.CCC.2023.29; CC BY4.0}
\begin{document}
\maketitle
\setcounter{section}{2}
\noindent\textit{Editorial note.} Original authors: Bruno P. Cavalar and Igor C. Oliveira. arXiv:2305.06821v1 (2023-05-11), https://arxiv.org/abs/2305.06821v1. Creative Commons Attribution4.0.\par\medskip
\noindent\textit{Editorial note.} License: https://creativecommons.org/licenses/by/4.0/\par\medskip
\noindent\textit{Editorial note.} Selected complete native author theorem/proof/context and full bibliography in standalone layout. Native input/mathhead and mathematical style definitions flattened as explicit spans; unused frontmatter and figure packages omitted. No mathematical changes or new proofs. Original CC BY4.0 license retained.\par\medskip
\noindent\textit{Editorial note.} Third-party diagram code: postlattice.sty, copyright2016 Gernot Salzer, maintained under LaTeX Project Public License1.3c or later, https://www.latex-project.org/lppl/ . Original package header/code preserved; mechanical flattening excludes only terminal endinput and wraps catcodes for standalone compilation. The mathematical paper remains CC BY4.0.\par\medskip
\noindent\textit{Allowed standard prerequisites (editorial).}\begin{itemize}
\item Known monotone lower bounds, logspace algorithms and threshold circuit bounds. Exact author prior theorem statements retained with citations. Full author depth-reduction/padding support proofs included.
\end{itemize}\medskip
\ifdefined\hypertarget\hypertarget{sourcepointer-1}{}\fi
\noindent\textit{Source reference.} Original reference 3: number-theory/sources/ccc-2023-29/native/main.tex, source lines 954--955. Classification overview compares to companion circuit-separation section; no premise from that section is required beyond separately retained exact prior theorem.\par\medskip
\ifdefined\hypertarget\hypertarget{sourcepointer-2}{}\fi
\noindent\textit{Source reference.} Original reference 3.2: number-theory/sources/ccc-2023-29/native/main.tex, source lines 1119--1120. Side discussion compares the delivered classification to a separate graph separation, not used in classification proof.\par\medskip
\section{Preliminaries}

\subsection{Notation}

\subparagraph{Boolean functions.}

We denote by $\monfns$ the set of all monotone Boolean functions.
We define $\poly = \set{n \mapsto n^C : C \in \bbn}$.
A Boolean function
    $f : \blt^{\binom{n}{2}} \to \blt$ is said to be a graph property
    if $f(G) = f(H)$ for any two isomorphic graphs $G$ and $H$.
Let $\calf = \set{f_n}_{n \in \bbn}$
    be a sequence of graph properties, where $f_n$ is defined over undirected graphs on
    $n$ vertices.
    We say that $\calf$ is
    \emph{preserved under homomorphisms} if,
    whenever there is a homomorphism from a graph $G$ to a graph $H$,
    we have $\calf(G) \leq \calf(H)$.
    We denote by $\HomPres$ 
    the set of all 
    graph properties which are preserved under homomorphisms.
    Note that $\HomPres \sseq \monfns$.

    \subparagraph{Boolean circuits.}

    We denote by $\AC^0_d[s]$ 
    the family of Boolean functions computed by
    size-$s$,
    depth-$d$
    Boolean circuits with unbounded fan-in $\set{\land,\lor}$-gates and input literals from $\{x_1, \overline{x_1}, \ldots, x_n, \overline{x_n}\}$.
    We write $\AC^0[s]$ as a shorthand for
    $\bigcup_{d = 1}^\infty \AC^0_d[s]$,
    and $\AC^0$ as a shorthand of $\AC^0[n^{O(1)}] = \AC^0[\poly]$.
    We will also refer to $\AC^0_d[\poly]$
    by $\AC^0_d$.
We write $\DNF[s]$ to denote the family of Boolean functions
    computed by size-$s$ DNFs, where size is measured by  number of terms.
    We write $\CNF[s]$ analogously.
    We write $\SIZE[s]$ to denote
    the family of Boolean functions
    computed by size-$s$ circuits.
    We write $\DEPTH[d]$ to denote
    the family of Boolean functions
    computed by fan-in 2 circuits of depth $d$.
We denote by $\AC^0[\xor]$ the family of Boolean functions computed by
polynomial-size $\AC^0$ circuits with unbounded fan-in $\xor$-gates.

We denote by $\L$ the 
family of
Boolean functions
computed by  logspace machines,
and by $\NL$ the family of Boolean functions
computed by polynomial-time nondeterministic logspace machines.
Moreover, we denote by $\pL$
the family of Boolean functions computed by
polynomial-time nondeterministic logspace machines with a \emph{parity}
acceptance condition (i.e.,
an input is accepted if the number of accepting paths is odd).

\subparagraph{Circuit complexity.}

Given a circuit class $\calc$, we write
    $\monclass{C}$ to denote the
    monotone version of $\calc$.
Given a function $f$, we write $\Cmon{f}$
    to denote the size of the smallest monotone circuit computing
    $f$ and $\mdepth{f}$ to denote the smallest depth of a fan-in 2
    monotone circuit computing $f$.
Given two Boolean functions $f,g$, we write
$f \mprojred g$ if
there exists a many-one reduction from $f$ to $g$
in which each bit of the reduction is a monotone
projection\footnote{A monotone projection is a projection without
negations.}
of the input.

\subparagraph{Miscellanea.}

Let $\alpha \in \blt^n$ .
    We 
    define
    $\hwt{\alpha} := \sum_{i=1}^n \alpha_i$.
    We call $\card{\alpha}_1$ the \defx{Hamming weight} of $\alpha$.
    We let
    $\supp(\alpha) = \set{i \in [n] : \alpha_i = 1}$.
We let $\thr_{k,n} : \blt^n \to \blt$ be the Boolean function
    such that
    $\thr_{k,n}(x) = 1 \iff \hwt{x} \geq k$.



\subsection{Background results}


The next lemma, which is proved via a standard ``guess-and-verify''
approach, shows that nondeterministic logspace computations can be
simulated by circuits of size $2^{n^{\varepsilon}}$ and of depth $d =
O_\varepsilon(1)$.

\begin{lemma}[Folklore; see, e.g.,~{\cite[Lemma 8.1]{AHMPS08}}]
    \label{lemma:ac0_nc1_sim}
    For all $\eps > 0$, we have
    $\NL \sseq \AC^0[2^{n^\eps}]$.
\end{lemma}







\begin{theorem}[\cite{gkrs19}]
    \label{thm:xor-sat}
    There exists $\eps > 0$ 
    such that 
    $\semMon{\pL} \not\sseq \mSIZE[2^{o(n^{\eps})}]$.
    Moreover, 
    this separation is witnessed by $\txorsat$.
\end{theorem}

\begin{theorem}[Polynomial-size constant-depth vs.~larger monotone depth]
    \label{thm:ac0-fml-graph-intro}
    For every $k \geq 1$,
    we have
    $\semMon{\AC^0} \not\sseq \mDEPTH[(\log n)^k]$. Moreover, this separation holds for a monotone graph property.
\end{theorem}

\begin{restatable}[Polynomial-size constant-depth vs.~larger monotone size]{theorem}{acxormp}
    \label{thm:ac02-mp-intro}
    For every $k \geq 1$,
    we have
$\semMon{\AC^0[\oplus]} \not\sseq \mSIZE[2^{(\log n)^k}]$.
\end{restatable}

\label{sec:intro_mon_csp}

Recall that 
\citep{gkrs19} 
showed the existence of a monotone function $f^{\mathsf{GKRS}}$ in $\NC^2$ that is not in $\mSIZE[2^{n^{\Omega(1)}}]$. As opposed to classical results \citep{razborov1985lower, andreev1985method, DBLP:journals/combinatorica/AlonB87, tardos_1988} that rely on the approximation method, their monotone circuit lower bound employs a  lifting technique from communication complexity. It is thus natural to consider if their approach can be adapted to provide a monotone function $g$ that is efficiently computable by constant-depth circuits but is not in $\mSIZE[\poly]$. 

As remarked in~\cite{gkrs19,DBLP:journals/sigact/RezendeGR22}, all
monotone lower bounds 
obtained from lifting theorems so far 
also hold for monotone encodings of constraint satisfaction problems (CSPs). 
Next, we introduce a class of monotone Boolean functions $\CSPf{S}$ which
capture the framework and lower bound of \citep{gkrs19}.\\

\noindent \textbf{Encoding CSPs as monotone Boolean functions.} Let $R \sseq \blt^k$ be a relation.
We call $k$ the \emph{arity} of $R$.
Let
$V = (i_1,\dots,i_k) \in [n]^k$,
and 
let
$f_{R,V}: \blt^n \to \blt$
be
the function 
that
accepts a string $x \in \blt^n$
if $(x_{i_1},\dots,x_{i_k}) \in R$.
We call $f_{R,V}$ a \emph{constraint application} of $R$ on $n$ variables.
(A different choice of the sequence $V$
gives a different constraint application of $R$.)
If $S$ is a finite set of Boolean relations,
we call any set of constraint applications of relations from $S$
on a fixed set of variables
an \emph{$S$-formula.}
In particular, we can describe an $S$-formula through a set of pairs $(V,R)$.
We say that an $S$-formula $F$ is \emph{satisfiable}
if there exists an assignment to the variables of $F$
which satisfies all the constraints of $F$.

Let $S = \set{R_1,\dots,R_k}$ be a finite set of Boolean relations.
Let $\ell_i$ be the arity of the relation $R_i$.
Note that there are $n^{\ell_i}$ possible constraint applications
of the relation $R_i$ on $n$ variables.
Let
$N := \sum_{i=1}^k n^{\ell_i}$.
We can identify each 
$S$-formula $F$
on a fixed set of $n$ variables 
with a corresponding string $w^F \in \blt^N$, where $w^F_j = 1$ if and only if the $j$-th possible constraint application (corresponding to one of the $N$ pairs $(V,R))$ appears in $F$.
Let $\CSPfn{S}{n} : \blt^N \to \blt$
be the Boolean function which accepts a given 
$S$-formula $F$ if $F$ is \emph{unsatisfiable}.
Note that this is a monotone function.
When $n$ is clear from the context or we view $\{\CSPfn{S}{n}\}_{n \geq 1}$ as a sequence of functions, we simply write
$\CSPf{S}$.\\

\begin{theorem}[Limits of the direct approach via lifting and CSPs]
    \label{thm:cspsat-intro}
    Let $S$ be a finite set of Boolean relations. The following holds. 
    \begin{enumerate}
        \item    If 
            $\CSPf{S} \notin \mSIZE[\mathsf{poly}]$
            then
            $\CSPf{S}$ is 
             $\pL$-hard under
            $\acmred$ reductions.
        \item  
            If
            $\CSPf{S} \notin \mNC^1$
            then $\CSPf{S}$ is $\L$-hard under
            $\acmred$ reductions.
    \end{enumerate}
\end{theorem}

\begin{theorem}
    \label{thm:ac0-fml-graph}
    For every $i \geq 1$,
    we have
    $\semMon{\AC^0} \cap \GraphPpts \not\sseq \mDEPTH[(\log n)^i]$.
    In particular, we have
    $\semMon{\AC^0} \cap \GraphPpts \not\sseq \mNC^i$.
\end{theorem}

\begin{theorem}[\cite{DBLP:journals/combinatorica/BabaiGW99}]
    \label{thm:bip-oddfactor-lb}
        Any monotone formula computing 
        $\OddFactor_n$ 
        has size $2^{\Omega(n)}$,
        and any monotone circuit computing
        $\OddFactor_n$ 
        has size $n^{\Omega(\log n)}$.
\end{theorem}

\begin{theorem}[Dichotomies for the monotone complexity of Boolean-valued CSPs]
    \label{thm:intro-mon-dichotomies}
    Let $S$ be a finite set of Boolean relations. The following holds.
    \begin{enumerate}
        \item \emph{Monotone Size Dichotomy:} If 
    $\Pol(S) \sseq \clonefont{L_3}$
     there is $\eps > 0$
    such that
    $\Cmon{\CSPf{S}} = 2^{\Omega(n^\eps)}$.
    Otherwise, 
    $\Cmon{\CSPf{S}} = n^{O(1)}$.

\item \emph{Monotone Depth Dichotomy:}
    If 
    $\Pol(S) \sseq \clonefont{L_3}$ or
    $\Pol(S) \sseq \clonefont{V_2}$ or
    $\Pol(S) \sseq \clonefont{E_2}$,
     there is $\eps > 0$
    such that
    $\mdepth{\CSPf{S}} = {\Omega(n^\eps)}$.
    Otherwise, 
    $\CSPf{S} \in \mNC^2$.
    \end{enumerate}
    
\end{theorem}

\section{Monotone Complexity of Constraint Satisfaction Problems}
\label{sec:CSPs}

In this section, we study the monotone complexity of Boolean-valued $\CSP$s.
Our goal is to classify which types of Boolean $\CSP$s are hard for
monotone circuit size and monotone circuit depth,
eventually proving
\Cref{thm:cspsat-intro,thm:intro-mon-dichotomies}.

We will first spend some time recalling standard definitions and concepts
in the theory of CSPs (\Cref{sec:csp-def}),
as well as a few results about CSPs that were
proved in previous works
\cite{DBLP:conf/stoc/Schaefer78,DBLP:journals/tcs/Jeavons98,
playing_one,playing_two,DBLP:journals/jcss/AllenderBISV09}
(\Cref{sec:csp-fact}).
We will then prove \Cref{thm:intro-mon-dichotomies}
in \Cref{sec:csp-dichotomy},
and
we will finally prove \Cref{thm:cspsat-intro}
in \Cref{sec:csp-conseq}
after 
proving some auxiliary results
in \Cref{sec:csp-ac0}.

\subsection{Definitions}
\label{sec:csp-def}

For a good introduction to the concepts defined below, we refer the reader
to~\cite{playing_one, playing_two}.
We also refer the reader to 
\Cref{sec:intro_mon_csp} 
for 
the definition of the family of functions
$\CSPf{S}$, as well as the terms
\emph{constraint application, $S$-formula} and \emph{satisfiable formula}.



We denote by $p_i^n : \blt^n \to \blt$ the
$i$-th \emph{projection function} on $n$ variables,
whose operation is defined as
$p_i^n(x) = x_i$.
For a set of Boolean functions $B$,
we
denote by $[B]$ the \emph{closure} of $B$,
defined as follows:
a Boolean function $f$ is in $[B]$
if and only if
$f \in B \cup \set{\clonefont{Identity}}$
or if there exists $g \in B$
and $h_1,\dots,h_k$
such that
$f = g(h_1,\dots,h_k)$,
where each $h_i$ is either a 
projection function or a function from $[B]$.
We can equivalently define $[B]$ as
the set of all Boolean functions
that can be computed by circuits using 
the functions of $B$ as gates.
Note that $[B]$ necessarily contains an infinite number of Boolean
functions, 
since $p_1^n \in [B]$ for every $n \in \bbn$;
moreover, the constant functions
are not necessarily in $[B]$.
We say that $B$ is a \emph{clone} if $B = [B]$.
A few prominent examples of clones are the set of all Boolean functions
(equal to $[\set{\land,\neg}]$),
monotone functions
(equal to $[\set{\land,\lor,0,1}]$), 
and linear functions
(equal to $[\set{\xor,1}]$).

\begin{remark}
    \label{rk:post_lattice}
    The set of all clones forms a lattice, known as
    \defx{Post's lattice},
    under the operations
    $[A] \sqcap [B] := [A] \cap [B]$
    and
    $[A] \sqcup [B] := [A \cup B]$.
    From the next section onwards,
    we will refer to the clones defined in~\cite{playing_one}
    (such as $\clonefont{I_0}$, $\clonefont{I_1}$, etc.),
    assuming the reader is familiar with them.
    For the unfamiliar reader, we refer to \Cref{sec:clone-background} and 
    \Cref{fig:clone_table,fig:post_lattice},
    which contain all the definitions of the clones we will need,
    as well as the entire Post's lattice in graphical representation.
    
    To avoid confusion, we will always refer to clones with
    $\clonefont{normal}$-$\clonefont{Roman}$ font (e.g.,
    $\clonefont{S_1}, \clonefont{I_0}$, etc).
\end{remark}

Let $S$ be a finite set of Boolean relations.
We denote by $\CNF(S)$ the set of all $S$-formulas.
We denote by $\COQ(S)$ 
the set of all relations which can be expressed
with the following type of formula $\varphi$:
\begin{equation*}
    \varphi(x_1,\dots,x_k) = \exists y_1,\dots,y_\ell
    \,
    \psi(x_1,\dots,x_k,y_1,\dots,y_\ell),
\end{equation*}
where $\psi \in \CNF(S)$.
The relations in $\COQ(S)$ will also be referred as \emph{conjunctive
queries} over $S$.
We denote by $\ccln{S}$ the set of relations
defined as
$\ccln{S} := \COQ(S \cup \set{=})$.
If $S = \ccln{S}$, we say that $S$ is a \emph{co-clone}.
We define 
$$\CSP = \set{\CSPf{S} : \text{$S$ is a finite set of relations}}.$$
We say that $\CSPf{S}$ is \emph{trivial}
if $\CSPf{S}$ is a constant function.


Let $R$ be a $k$-ary Boolean relation 
and
let $f : \blt^\ell \to \blt$ be a Boolean function.
For $x \in R$ and $i \in [k]$, we denote by $x[i]$
the $i$-th bit of $x$.
\begin{definition}
    \label{def:polymorphism}
    We say that $f$ is a
    \defx{polymorphism} of $R$,
    and $R$ is 
    \defx{an invariant of $f$},
    if, for all $x_1, \dots, x_\ell \in R$, we have
    \begin{equation*}
        (
            f(x_1[1],\dots,x_{\ell}[1]),
            f(x_2[2],\dots,x_{\ell}[2]),
            \dots,
            f(x_k[k],\dots,x_{\ell}[k])
        )
        \in
        R.
    \end{equation*}
    We denote the set of all polymorphisms of $R$ by $\Pol(R)$.
    For a set of relations $S$,
    we denote by $\Pol(S)$ the set of Boolean functions
    which are polymorphisms of all the relations of $S$.
    For a set of Boolean functions, we denote by
    $\Inv(B)$ the set of all Boolean relations which are invariant
    under all functions of $B$ (i.e.,
    $\Inv(B) = \set{R : B \sseq \Pol(R)}$).
\end{definition}

The following summarises the important facts about
clones, co-clones and polymorphisms that are relevant to
the study of CSPs~\cite{10.1145/263867.263489}. 

\begin{lemma}
    \label{lem:galois}
    Let 
    $S$ 
    and $S'$
    be sets
    of Boolean relations 
    and
    let
    $B$ 
    and $B'$
    be sets
    of Boolean functions.
    We have
    \begin{enumerate}[label=(\roman*)]
        \item 
            \label{item:pol-clone}
            $\Pol(S)$ is a clone and $\Inv(B)$ is a co-clone;
        \item 
            \label{item:pol-invert}
            If 
            $S \sseq S'$,
            then
            $\Pol(S') \sseq \Pol(S)$;
        \item 
            \label{item:inv-invert}
            If 
            $B \sseq B'$,
            then
            $\Inv(B') \sseq \Inv(B)$;
        \item 
            \label{item:coq-idempotent}
            $\COQ(\COQ(S)) = \COQ(S)$;
        \item 
            \label{item:coq-sseq}
            If $S \sseq S'$, then
            $\COQ(S) \sseq \COQ(S')$;
        \item 
            \label{item:invpol}
            $\Inv(\Pol(S)) = \ccln{S}$;
        \item 
            \label{item:polinv}
            $\Pol(\Inv(B)) = [B]$.
    \end{enumerate}
\end{lemma}

We now define 
different
types of reductions.
We say that a reduction is a \emph{monotone OR-reduction}
if
every bit of the reduction is either constant
or can be computed by a monotone disjunction on the
input variables.
We write $f \mormred g$ if 
there exists a many-one monotone OR-reduction
from $f$ to $g$.
We also write 
$f \acmred g$
if 
there exists a many-one
$\AC^0$ reduction
from $f$ to $g$,
and
$f \mnlmred g$
if 
there exists a many-one
$\mNL$ reduction
from $f$ to $g$\footnote{A many-one $\AC^0$ (resp. $\mNL$) reduction is one
in which each bit of the reduction is either constant or 
can be computed with a polynomial-size $\AC^0$ circuit
(resp. monotone nondeterministic branching program). 
Recall that a 
\emph{monotone nondeterministic branching program} 
is a directed acyclic graph $G$ with two
distinguished vertices $s$ and $t$, 
in which each edge $e$ is labelled with an input function 
$\rho_e \in \set{1,x_1,\dots,x_n}$. 
Given an input $x$, the program accepts if there
exists a path from $s$ to $t$ in the subgraph $G_x$ of $G$ 
in which an edge $e$ appears if $\rho_e(x)=1$.
}.
Unless otherwise specified, every reduction we consider will generate an instance
of polynomial size on the length of the input.

Finally, we denote by $\OR^k$ and $\NAND^k$
the $k$-ary $\OR$ and $\NAND$ relations, respectively.


\subsection{Basic facts about $\CSPSAT$}
\label{sec:csp-fact}



We state 
here basic facts about the $\CSPSAT$ function.
These facts are proved in the original paper of
Schaefer~\cite{DBLP:conf/stoc/Schaefer78},
as well as in later
papers~\cite{DBLP:journals/tcs/Jeavons98,
playing_one,playing_two,DBLP:journals/jcss/AllenderBISV09}.

\Cref{prop:poly} below is one of the most important lemmas of this section and
will be used many times. It states that $\Pol(S)$
characterises the monotone complexity of $\CSPf{S}$, in the sense that 
the sets of relations with few polymorphisms give rise to the hardest
instances of CSPs.
A non-monotone version of this result was proved in
\cite[Theorem 2.4]{DBLP:journals/tcs/Jeavons98, playing_two},
and we
check 
in \Cref{sec:reduction_monotone}
that their proofs also hold in the monotone
case.





 

\begin{restatable}[\protect{Polymorphisms characterise the complexity of
    CSPs~\cite[Theorem 2.4]{DBLP:journals/tcs/Jeavons98, playing_two}}]
    {lemma}{polym}
    \label{prop:poly}
    If $\Pol(S_2) \sseq \Pol(S_1)$,
    then
    $\CSPfn{S_1}{n} \mnlmred \CSPfn{S_2}{\poly(n)}$.
\end{restatable}




\Cref{thm:mon-schaefer} gives monotone circuit upper bounds
for some instances of $\CSPf{S}$.
Non-monotone variants of this upper bound were originally obtained
in the seminal paper of Schaefer~\cite{DBLP:conf/stoc/Schaefer78},
and we again check that the monotone variants work in
\Cref{sec:monotone_upper}.

\begin{restatable}[Monotone version of the upper bounds 
    for $\CSPSAT$ 
\cite{DBLP:conf/stoc/Schaefer78,DBLP:journals/jcss/AllenderBISV09}]
{theorem}
{monschaeffer}
    \label{thm:mon-schaefer}
    Let $S$ be a finite set of relations.
    The following holds.
    \begin{enumerate}
        \item
            If $\clonefont{\clonefont{E_2}} \sseq \Pol(S)$ or $\clonefont{\clonefont{V_2}} \sseq \Pol(S)$,
            then
            $\CSPf{S} \in \mSIZE[\poly]$.
        \item 
            If 
            $\clonefont{\clonefont{D_2}} \sseq \Pol(S)$,
            or
            $\clonefont{\clonefont{S_{00}}} \sseq \Pol(S)$,
            or
            $\clonefont{\clonefont{S_{10}}} \sseq \Pol(S)$,
            then
            $\CSPf{S} \in \mNL$.
    \end{enumerate}
\end{restatable}

Finally, 
we state here a result of~\cite{DBLP:journals/jcss/AllenderBISV09},
which
classifies the \emph{non-monotone} complexity of $\CSPf{S}$
under $\acmred$ reductions.
The classification of the complexity of $\CSPf{S}$
is based solely on $\Pol(S)$.
See \Cref{fig:post_lattice} for a graphical representation.

\begin{theorem}[\protect{Refined classification of $\CSP$ problems 
    \cite[Theorem 3.1]{DBLP:journals/jcss/AllenderBISV09}}]
    \label{thm:refined-completeness}
    Let $S$ be a finite set of Boolean relations.
    The following holds.
    \begin{itemize}
        \item If 
            $\clonefont{I_0} \sseq \Pol(S)$
            or 
            $\clonefont{I_1} \sseq \Pol(S)$,
            then
            $\CSPf{S}$ is trivial.
        \item 
            If 
            $\Pol(S) \in \set{\clonefont{I_2},\clonefont{N_2}}$,
            then
            $\CSPf{S}$
            is
            \acmcmpl{}
            for
            $\NP$.
        \item 
            If 
            $\Pol(S) \in 
            \set{\clonefont{V_2},\clonefont{E_2}}$,
            then
            $\CSPf{S}$
            is
            \acmcmpl{}
            for
            $\P$.
        \item 
            If 
            $\Pol(S) \in 
            \set{\clonefont{L_2},\clonefont{L_3}}$,
            then
            $\CSPf{S}$
            is
            \acmcmpl{}
            for
            $\pL$.
        \item 
            If 
            $\clonefont{S_{00}} \sseq \Pol(S) \sseq \clonefont{S_{00}}^2$
            or
            $\clonefont{S_{10}} \sseq \Pol(S) \sseq \clonefont{S_{10}}^2$
            or
            $\Pol(S) \in 
            \set{\clonefont{D_2},\clonefont{M_2}}$,
            then
            $\CSPf{S}$
            is
            \acmcmpl{}
            for
            $\NL$.
        \item 
            If 
            $\Pol(S) \in 
            \set{\clonefont{D_1}, \clonefont{D}}$,
            then
            $\CSPf{S}$
            is
            \acmcmpl{}
            for
            $\L$.
        \item 
            If 
            $\clonefont{S_{02}} \sseq \Pol(S) \sseq \clonefont{R_2}$
            or
            $\clonefont{S_{12}} \sseq \Pol(S) \sseq \clonefont{R_2}$,
            then
            either
            $\CSPf{S} \in \AC^0$
            or
            $\CSPf{S}$
            is
            \acmcmpl{}
            for
            $\L$.
    \end{itemize}
\end{theorem}

\newcommand\ABISV
    {R2/.cstyle=LN,
     M/.cstyle=NA,
     M2/.cstyle=NL,
     D/.cstyle=L,
     D2/.cstyle=NL,
     S212/.cstyle=LN,
     S312/.cstyle=LN,
     S12/.cstyle=LN,
     S202/.cstyle=LN,
     S302/.cstyle=LN,
     S02/.cstyle=LN,
     S00/.cstyle=NL,
     S300/.cstyle=NL,
     S200/.cstyle=NL,
     S10/.cstyle=NL,
     S310/.cstyle=NL,
     S210/.cstyle=NL,
     E2/.cstyle=P,
     V2/.cstyle=P,
     L2/.cstyle=pL,
     L3/.cstyle=pL,
     N2/.cstyle=NP,
     I2/.cstyle=NP,
     xscale=-1, yscale=-1,scale=1,labels=wagner,cshape/circle/.append style={minimum width=15pt}
    }
\newlength{\boxwidth}
\settowidth{\boxwidth}{$\NP$-complete}
\addtolength{\boxwidth}{5mm}
\begin{figure}
    \centering
    \expandafter\postlattice\expandafter[\ABISV]
    \bigskip

    \begin{tabular}{lll}
        \makebox[\boxwidth][l]{\protect\postlegend{NP} $\NP$-complete}
        & \makebox[\boxwidth][l]{\protect\postlegend{P} $\P$-complete}
        & \makebox[\boxwidth][l]{\protect\postlegend{pL} $\pL$-complete}
        \\
        \makebox[\boxwidth][l]{\protect\postlegend{NL} $\NL$-complete}
        & \makebox[\boxwidth][l]{\protect\postlegend{L} $\L$-complete}
        & \makebox[\boxwidth][l]{\protect\postlegend{LN} $\L$-complete or in  $\mathsf{coNLOGTIME}$}
        \\
        & \makebox[\boxwidth][l]{\protect\postlegend{NA} Trivial}
        &
        \\
    \end{tabular}

    \caption{
        Graph of all closed classes of Boolean functions.
        The vertices are colored with the complexity of 
        deciding $\CSP$s 
        whose set of polymorphisms corresponds to the label of the vertex.
        Trivial $\CSP$s are those that correspond to constant functions.
        Every hardness result is proved under $\acmred$ reductions.
        See \Cref{thm:refined-completeness} for details.
        A similar figure appears in
        \cite[Figure 1]{DBLP:journals/jcss/AllenderBISV09}.}
    \label{fig:post_lattice}
\end{figure}



\subsection{A monotone dichotomy for $\CSPSAT$}
\label{sec:csp-dichotomy}

In this section, we prove \Cref{thm:intro-mon-dichotomies}.
We first prove Part (1) of the theorem (the dichotomy for circuit size),
and then we prove Part (2) of the theorem (the dichotomy for circuit
depth).

\subparagraph{Dichotomy for circuits.}

To prove the dichotomy for circuits, 
we first show that, for
any set of relations $S$ whose set of polymorphisms is contained in $\clonefont{L_3}$,
we can monotonically reduce $\txorsat$ to $\CSPf{S}$.


\begin{lemma}
    \label{lem:l3_xorsat}
    Let $S$ be a finite set of relations.
    If 
    $\Pol(S) \sseq \clonefont{L_3}$,
    then
    $\txorsat \mnlmred \CSPf{S}$.
\end{lemma}
\begin{proof}
    Inspecting Post's lattice
    (\Cref{fig:post_lattice}),
    note that the only clones 
    strictly contained in
    $\clonefont{L_3}$ are
    $\clonefont{L_2}, \clonefont{N_2}$ and $\clonefont{I_2}$.
    We will first show that the reduction holds for the case
    $\Pol(S)=\clonefont{L_2}$ and then prove that 
    the reduction also holds for the case $\Pol(S)=\clonefont{L_3}$.
    \Cref{prop:poly}
    will then imply the cases $\Pol(S) \in \set{\clonefont{N_2},\clonefont{I_2}}$, since
    $\clonefont{I_2} \sseq \clonefont{N_2} \sseq \clonefont{L_3}$.


    It's not hard to check that,
    if $\Pol(S) = \clonefont{L_2}$,
    then $\Pol(S) \sseq \Pol(\txorsat)$
    (it suffices to observe that bitwise XORing
    three satisfying assignments to a linear equation gives rise to 
    a new satisfying assignment to the same equation).
    Therefore, from \Cref{prop:poly} we deduce
    that $\txorsat$ admits a reduction to
    $\CSPf{S}$ in $\mNL$.
    In order to prove the case $\Pol(S) = \clonefont{L_3}$, we first prove the
    following claim.

    \begin{claim*}[\protect{\cite[Lemma 3.11]{DBLP:journals/jcss/AllenderBISV09}}]
        \label{lem:l2-to-l3}
        Let $S$ be a finite set of relations such that
        $\Pol(S) = \clonefont{L_2}$.
        There exists a finite set of relations $S'$
        such that
        $\Pol(S') = \clonefont{L_3}$
        and
        $\CSPfn{S}{n} \mprojred \CSPfn{S'}{n+1}$.
    \end{claim*}
    \begin{proof}
        We describe the proof of Lemma 3.11
        in~\cite{DBLP:journals/jcss/AllenderBISV09}
        and observe that it gives a monotone reduction.

        For a relation $R \in S$,
        let $R' = \set{(\neg x_1,\dots, \neg x_k) : (x_1,\dots,x_k) \in R}$.
        Let also $S' = \set{R' : R \in S}$.
        It's not hard to check that 
        $\Pol(S') = \clonefont{L_3}$,
        since $S'$ is 
        an invariant of
        $\clonefont{L_2}$ and $\clonefont{N_2}$,
        and $\clonefont{L_3}$
        is the smallest clone containing both $\clonefont{L_2}$ and $\clonefont{N_2}$;
        moreover, if $\rho \in \Pol(S')$ and $\rho$ is a Boolean function on at
        least two bits, then
        $\rho \in \Pol(S) = \clonefont{L_2}$.

        Now let 
        $F$
        be a instance of $\CSPfn{S}{n}$.
        For every constraint $C = R(x_1,\dots,x_k)$ in $F$,
        we
        add the constraint
        $C' = R'(\alpha,x_1,\dots,x_k)$ to the $S'$-formula $F'$,
        where $\alpha$ is a new variable.
        Note that $F'$ is a $S'$-formula, defined on $n+1$ variables,
        which
        is satisfiable 
        if and only if $F$ is satisfiable.
        Moreover, the construction of $F'$ from $F$ can be done with a monotone
        projection.
    \end{proof}

    Since the case $\Pol(S) = \clonefont{L_2}$ holds,
    the case $\Pol(S) = \clonefont{L_3}$ now follows
    from~\Cref{prop:poly} and the Claim.
    Finally, from \Cref{prop:poly}
    we conclude that the 
    reduction
    also holds for the case
    $\Pol(S) \in \set{\clonefont{N_2},\clonefont{I_2}}$, since
    $\clonefont{I_2} \sseq \clonefont{N_2} \sseq \clonefont{L_3}$.
\end{proof}


\begin{theorem}[Dichotomy for monotone circuits]
    \label{thm:monckt-dichotomy}
    Let $S$ be a finite set of relations.
    If 
    $\Pol(S) \sseq \clonefont{L_3}$
    then there exists a constant $\eps > 0$
    such that
    $\Cmon{\CSPf{S}} = 2^{\Omega(n^\eps)}$.
    Otherwise, 
    we have
    $\Cmon{\CSPf{S}} = n^{O(1)}$.
\end{theorem}
\begin{proof}
    If $\Pol(S) \sseq \clonefont{L_3}$, the lower bound follows from
    the 'moreover' part of 
    \Cref{thm:xor-sat}, and 
    \Cref{lem:l3_xorsat}.
    For the upper bound, we inspect Post's
    lattice~(\Cref{fig:post_lattice}).
    Observe that, if $\Pol(S) \not\sseq \clonefont{L_3}$, the following are the only possible cases:
    \begin{enumerate}
        \item $\clonefont{I_0} \sseq \Pol(S)$ or $\clonefont{I_1} \sseq \Pol(S)$.
            In both cases, any $\CNF(S)$ is trivially satisfiable.
        \item $\clonefont{E_2} \sseq \Pol(S)$ or $\clonefont{V_2} \sseq \Pol(S)$.
            In this case, $\CSPf{S} \in \mSIZE[\poly]$ 
            by~\Cref{thm:mon-schaefer}.
        \item $\clonefont{D_2} \sseq \Pol(S)$.
            In this case, $\CSPf{S} \in \mNL \sseq \mSIZE[\poly]$ 
            by~\Cref{thm:mon-schaefer}.
            \qedhere
    \end{enumerate}
\end{proof}

\tikzset{NA/.style  = {fill = white}}
\tikzset{mP/.style  = {fill = magenta!50}}
\tikzset{LB/.style  = {pattern = north east lines, pattern color = applegreen!80, line width=1.5pt}}
\newcommand\CktDichotomy
    {R2/.cstyle=mP,
     M/.cstyle=NA,
     M2/.cstyle=mP,
     D/.cstyle=mP,
     D2/.cstyle=mP,
     S212/.cstyle=mP,
     S312/.cstyle=mP,
     S12/.cstyle=mP,
     S202/.cstyle=mP,
     S302/.cstyle=mP,
     S02/.cstyle=mP,
     S00/.cstyle=mP,
     S300/.cstyle=mP,
     S200/.cstyle=mP,
     S10/.cstyle=mP,
     S310/.cstyle=mP,
     S210/.cstyle=mP,
     E2/.cstyle=mP,
     V2/.cstyle=mP,
     L2/.cstyle=LB,
     L3/.cstyle=LB,
     N2/.cstyle=LB,
     I2/.cstyle=LB,
     xscale=-1, yscale=-1,scale=1,labels=wagner,cshape/circle/.append style={minimum width=15pt}
    }
\newlength{\cktdichotomybox}
\settowidth{\cktdichotomybox}{Requires monotone circuits of large size}
\addtolength{\cktdichotomybox}{5mm}
\begin{figure}
    \centering
    \expandafter\postlattice\expandafter[\CktDichotomy]
    \bigskip

    \begin{tabular}{l}
        \makebox[\cktdichotomybox][l]{\protect\postlegend{mP} Solvable in
        $\mSIZE[\poly]$.}
        \\
        \makebox[\cktdichotomybox][l]{\protect\postlegend{LB} Requires
        monotone circuits of size $2^{\Omega(n^{\eps})}$.}
        \\
        \makebox[\cktdichotomybox][l]{\protect\postlegend{NA} 
        Trivial (i.e., a constant function).}
    \end{tabular}

    \caption{Illustration of \Cref{thm:monckt-dichotomy}.
    The vertices are colored with the 
    \emph{monotone circuit size} complexity of deciding CSPs whose set of
    polymorphisms corresponds to the label of the vertex.
    }
    \label{fig:monckt-dichotomy}
\end{figure}



\begin{remark}
    \label{rk:lifting-approximation}
    We remark that the lifting theorem of~\cite{gkrs19}
    (which is an ingredient in the proof of \Cref{thm:xor-sat})
    is only used to prove that the monotone complexity of 
    $\CSPf{S}$
    is exponential
    when
    $\Pol(S) \sseq \clonefont{L_3}$.
    If we only care to show a superpolynomial separation, then 
    it suffices to apply the superpolynomial lower bound for CSPs 
    with counting proved
    in~\cite{DBLP:journals/siamcomp/FederV98,DBLP:journals/combinatorica/BabaiGW99}
    using the approximation method.
    Indeed, we give an explicit proof in~\Cref{sec:xorsat-appr}.
    The same holds for the consequences of this theorem
    (see \Cref{thm:cspsat-hard}).
\end{remark}

\subparagraph{Dichotomy for formulas.}


Define $\thornsat_n : \blt^{2n^3+n} \to \blt$ 
as $\thornsat_n = \CSPfn{\hornt}{n}$, where
$$
\hornt = \set{
    (\neg x_1 \lor \neg x_2 \lor x_3),
    (\neg x_1 \lor \neg x_2 \lor \neg x_3),
    (x)
}.$$
The following is proved 
in~\cite{DBLP:journals/combinatorica/RazM99, gkrs19}.


\begin{theorem}[\cite{DBLP:journals/combinatorica/RazM99, gkrs19}]
    \label{thm:3hornsat_lb}
    There exists $\eps > 0$
    such that
    $\thornsat \in \mSIZE[\poly] \sm \mDEPTH[o(n^\eps)]$.
\end{theorem}
\begin{proof}[Proof sketch]
    Since $\clonefont{E_2} \sseq \Pol(\hornt)$ (see,
    e.g.,~\cite[Lemma 4.8]{DBLP:books/daglib/0004131}),
    the upper bound follows from~\Cref{thm:mon-schaefer}.
    The lower bound follows from a lifting theorem
    of~\cite{DBLP:journals/combinatorica/RazM99,gkrs19}.
    They show that 
    the monotone circuit-depth of $\thornsat$
    is at least the
    depth of the smallest Resolution-tree
    refuting a so-called \emph{pebbling formula}.
    Since this formula requires Resolution-trees of depth $n^\eps$,
    the lower bound follows.
\end{proof}

Analogously to the previous section, we show that $\thornsat$ reduces
to $\CSPf{S}$ whenever
$\Pol(S)$ is small enough, in 
a precise sense stated below.
We then deduce the dichotomy for formulas with a similar argument.

\begin{lemma}
    \label{lem:l3_hornsat}
    Let $S$ be a finite set of relations.
    If 
    $\Pol(S) \sseq \clonefont{E_2}$ or $\Pol(S) \sseq \clonefont{V_2}$,
    then
    $\thornsat \mnlmred \CSPf{S}$.
\end{lemma}
\begin{proof}
    We first consider the case $\Pol(S) \sseq \clonefont{E_2}$.
    Note that
    $\clonefont{E_2} \sseq \Pol(\thornsat)$
    (see, e.g., \cite[Lemma 4.8]{DBLP:books/daglib/0004131}).
    Therefore, from \Cref{prop:poly} we deduce
    that $\thornsat$ admits a reduction to
    $\CSPf{S}$ in $\mNL$.

    Now let
    $\tahornsat = \CSPf{\ahornt}$,
    where
    $\ahornt = 
    \set{
        (x_1 \lor x_2 \lor \neg x_3),
        (x_1 \lor x_2 \lor x_3),
        (\neg x)
    }$.
    Observe that a
    $\hornt$-formula $\varphi$
    is satisfiable if and only if
    the $\ahornt$-formula 
    $\varphi(\neg x_1,\dots,\neg x_n)$
    is satisfiable.
    Therefore, we have
    $\thornsat \mprojred \tahornsat$.
    Observing that
    $\clonefont{V_2} \sseq \Pol(\ahornt)$
    (again, see e.g. \cite[Lemma 4.8]{DBLP:books/daglib/0004131}),
    the result now follows from
    \Cref{prop:poly} and the previous paragraph.
\end{proof}


\begin{theorem}[Dichotomy for monotone formulas]
    \label{thm:monfml-dichotomy}
    Let $S$ be a finite set of relations.
    If 
    $\Pol(S) \sseq \clonefont{L_3}$, or
    $\Pol(S) \sseq \clonefont{V_2}$, or
    $\Pol(S) \sseq \clonefont{E_2}$,
    then there 
    is
    a constant $\eps > 0$
    such that
    $\mdepth{\CSPf{S}} = {\Omega(n^\eps)}$.
    Otherwise, 
    we have
    $\CSPf{S} \in \mNL
    \sseq \mNC^2 \sseq \mDEPTH[\log^2 n]$.
\end{theorem}
\begin{proof}
    We will first prove the lower bound.
    If $\Pol(S) \sseq \clonefont{L_3}$, the lower bound follows
    from \Cref{thm:monckt-dichotomy}.
    If
    $\Pol(S) \sseq \clonefont{V_2}$ or
    $\Pol(S) \sseq \clonefont{E_2}$,
    the lower bound follows 
    from \Cref{thm:3hornsat_lb,lem:l3_hornsat}.
    
    By inspecting Post's 
    lattice~(\Cref{rk:post_lattice}),
    we see that
    the remaning cases are:
    \begin{enumerate}
        \item $\clonefont{I_0} \sseq \Pol(S)$ or $\clonefont{I_1} \sseq \Pol(S)$.
            In both cases, any $\CNF(S)$ is trivially satisfiable.
        \item 
            $\clonefont{S_{00}} \sseq \Pol(S)$,
            or
            $\clonefont{S_{10}} \sseq \Pol(S)$,
            or
            $\clonefont{D_2} \sseq \Pol(S)$.
            In all of those three cases, we have
            $\CSPf{S} \in \mNL$ by
            \Cref{thm:mon-schaefer}.
            \qedhere
    \end{enumerate}
\end{proof}

\tikzset{mD/.style = {fill = mikadoyellow}}
\tikzset{dLB/.style = {pattern = north east lines, pattern color = cyan(process)!70, line width=1.5pt}}
\newcommand\FmlDichotomy
    {R2/.cstyle=mD,
     M/.cstyle=NA,
     M2/.cstyle=mD,
     D/.cstyle=mD,
     D2/.cstyle=mD,
     S212/.cstyle=mD,
     S312/.cstyle=mD,
     S12/.cstyle=mD,
     S202/.cstyle=mD,
     S302/.cstyle=mD,
     S02/.cstyle=mD,
     S00/.cstyle=mD,
     S300/.cstyle=mD,
     S200/.cstyle=mD,
     S10/.cstyle=mD,
     S310/.cstyle=mD,
     S210/.cstyle=mD,
     E2/.cstyle=dLB,
     V2/.cstyle=dLB,
     L2/.cstyle=dLB,
     L3/.cstyle=dLB,
     N2/.cstyle=dLB,
     I2/.cstyle=dLB,
     xscale=-1, yscale=-1,scale=1,labels=wagner,cshape/circle/.append style={minimum width=15pt}
    }
\begin{figure}
    \centering
    \expandafter\postlattice\expandafter[\FmlDichotomy]
    \bigskip

    \begin{tabular}{l}
        \makebox[\cktdichotomybox][l]{\protect\postlegend{mD} Solvable in
        $\mNL \sseq \mDEPTH[O(\log^2 n)]$.}
        \\
        \makebox[\cktdichotomybox][l]{\protect\postlegend{dLB} Requires
        monotone depth ${\Omega(n^{\eps})}$.}
        \\
        \makebox[\cktdichotomybox][l]{\protect\postlegend{NA} 
        Trivial
        (i.e., a constant function).}
    \end{tabular}

    \caption{Illustration of \Cref{thm:monfml-dichotomy}.
        The vertices are colored with the 
        \emph{monotone circuit depth} complexity of deciding CSPs whose set of
        polymorphisms corresponds to the label of the vertex.
    }
    \label{fig:monfml-dichotomy}
\end{figure}



\subsection{Some auxiliary results}
\label{sec:csp-ac0}

In this section, we prove auxiliary results needed in the proof
of a more general form of \Cref{thm:cspsat-intro}.
In particular, we will prove that all $\CSPf{S}$ which are in $\AC^0$ 
are also contained in $\mAC^0 \sseq \mNC^1$.
Moreover, we show that,
if $\CSPf{S} \notin \mNC^1$,
then $\CSPf{S}$ is $\L$-hard under $\acmred$ reductions.

We first observe that, when $\COQ(S_1) \sseq \COQ(S_2)$,
there exists an efficient low-depth reduction
from $\CSPf{S_1}$ to $\CSPf{S_2}$.
This reduction, which will be useful in this section,
is more refined than the one given by
\Cref{prop:poly}.
A proof of the non-monotone version of this statement
is found in
\cite[Proposition 2.3]{playing_two},
and we give a monotone version of this proof in
\Cref{sec:reduction_monotone}.

\begin{restatable}[\protect{\cite[Proposition 2.3]{playing_two}}]{lemma}{inclcoq}
    \label{prop:incl-coq}
    If $\COQ(S_1) \sseq \COQ(S_2)$,
    then
    there exists a constant $C \in \bbn$ such that
    $\CSPfn{S_1}{n} \mormred \CSPfn{S_2}{C n}$.
\end{restatable}
\begin{proof}
    We defer the proof to
    \Cref{sec:reduction_monotone}.
\end{proof}

We now recall some lemmas
from~\cite{DBLP:journals/jcss/AllenderBISV09},
and prove a few consequences from them.
We say that a set $S$ of relations \emph{can express equality} if
$\set{=} \sseq \COQ(S)$.

\begin{lemma}[\cite{DBLP:journals/jcss/AllenderBISV09}]
    \label{lem:coq-no-eq}
    Let $S$ be a finite set of relations.
    Suppose
    $\clonefont{S_{02}} \sseq \Pol(S)$
    ($\clonefont{S_{12}} \sseq \Pol(S)$, resp.)
    and that $S$ cannot express equality.
    Then there exists $k \geq 2$
    such that
    $S \sseq \COQ(\set{\OR^k, x, \neg x})$
    ($S \sseq \COQ(\set{\NAND^k, x, \neg x})$, resp.).
\end{lemma}
\begin{proof}
    Follows from the proof of Lemma 3.8
    of~\cite{DBLP:journals/jcss/AllenderBISV09}.
\end{proof}

\begin{lemma}
    \label{lem:ac0-cspsat}
    Let $S$ be a finite set of relations
    such that $\Pol(S) \sseq \clonefont{R_2}$.
    If $\clonefont{S_{02}} \sseq \Pol(S)$
    or
    $\clonefont{S_{12}} \sseq \Pol(S)$,
    and $S$ cannot express equality,
    then
    $\CSPf{S} \in 
    \mAC^0_3$.
\end{lemma}
\begin{proof}
    We write the proof in the case
    $\clonefont{S_{02}} \sseq \Pol(S)$.
    The other case is analogous.

    From \Cref{prop:incl-coq,lem:coq-no-eq} and
    \Cref{item:coq-idempotent,item:coq-sseq} of \Cref{lem:galois},
    we get that
    there is a monotone OR-reduction
    from
    $\CSPf{S}$ to
    $\CSPf{\set{\OR^k, x, \neg x}}$ for some $k$.
    However, an $\set{\OR^k, x, \neg x}$-formula
    is unsatisfiable iff
    there exists a literal and its negation as a constraint in the formula,
    or if there exists a disjunction
    in the formula
    such that
    every one of its literals appears negatively as a constraint.
    This condition can be easily checked by a polynomial-size
    monotone DNF.
    Composing the monotone DNF with the monotone OR-reduction,
    we obtain a depth-3 $\AC^0$ circuit computing
    $\CSPf{S}$.
\end{proof}

\begin{lemma}[\protect{\cite[Lemma 3.8]{DBLP:journals/jcss/AllenderBISV09}}]
    \label{lem:equality-lhard}
    Let $S$ be a finite set of relations
    such that $\Pol(S) \sseq \clonefont{R_2}$.
    If 
    $\clonefont{S_{02}} \sseq \Pol(S)$
    or
    $\clonefont{S_{12}} \sseq \Pol(S)$,
    and $S$ can express equality,
    then
    $\CSPf{S}$ is $\L$-hard
    under $\acmred$ reductions.
\end{lemma}



\begin{lemma}
    \label{lem:everything-lhard}
    Let $S$ be a finite set of relations.
    If 
    $\clonefont{S_{02}} \not\sseq \Pol(S)$
    and
    $\clonefont{S_{12}} \not \sseq \Pol(S)$,
    then
    $\CSPf{S}$ is $\L$-hard or trivial.
\end{lemma}
\begin{proof}
    This
    follows by 
    inspecting
    Post's 
    lattice~(\Cref{fig:post_lattice})
    and the classification theorem~(\Cref{thm:refined-completeness}).
\end{proof}

We may now prove the main result of this subsection.

\begin{theorem}
    \label{thm:ac0-csp}
    We have
    $\CSP \cap \AC^0 \sseq \mAC^0_3$.
    Moreover, if $\CSPf{S} \notin 
    \mAC^0_3$,
    then
    $\CSPf{S}$ is $\L$-hard under $\acmred$ reductions.
\end{theorem}
\begin{proof}
    Let $S$ be a finite set of relations.
    If
    $\CSPf{S} \not\in 
    \mAC^0_3$,
    then,
    by \Cref{lem:ac0-cspsat},
    at least 
    one of the following cases hold:
    \begin{enumerate}
        \item $\clonefont{S_{02}} \sseq \Pol(S) \sseq \clonefont{R_2}$ 
            or 
            $\clonefont{S_{12}} \sseq \Pol(S) \sseq \clonefont{R_2}$,
            and $S$ can express the equality relation;
        \item 
            $\clonefont{S_{02}} \not\sseq \Pol(S) \sseq \clonefont{R_2}$
            and
            $\clonefont{S_{12}} \not \sseq \Pol(S) \sseq \clonefont{R_2}$.
        \item 
            $\Pol(S) \not\sseq \clonefont{R_2}$.
    \end{enumerate}
    Since $\CSPf{S}$ is not trivial,
    we obtain that $\CSPf{S}$ is $\L$-hard 
    in the first two cases
    by
    \Cref{lem:equality-lhard,lem:everything-lhard},
    and it's easy to check that $\CSPf{S}$ is also
    $\L$-hard in the third case
    by 
    inspecting
    Post's 
    lattice~(\Cref{fig:post_lattice})
    and the classification theorem~(\Cref{thm:refined-completeness}).
    Since $\L \not\sseq \AC^0$,
    this also implies that, if $\CSPf{S} \in \AC^0$,
    then
    $\clonefont{S_{02}} \sseq \Pol(S) \sseq \clonefont{R_2}$
    or
    $\clonefont{S_{12}} \sseq \Pol(S) \sseq \clonefont{R_2}$,
    and $S$ cannot express the equality relation.
    \Cref{lem:ac0-cspsat}
    again gives
    $\CSPf{S} \in \mAC^0_3$.
\end{proof}


\subsection{Consequences for monotone circuit lower bounds via lifting}
\label{sec:csp-conseq}

We now prove a stronger form of \Cref{thm:cspsat-intro}.
In the previous section, we showed that
$\CSP \cap \AC^0 \sseq \mAC^0$.
In particular, this means that there does not exist a finite set of relations $S$
such that
$\CSPf{S}$ separates $\AC^0$ and $\mNC^1$,
a separation which we proved in \Cref{thm:ac0-fml-graph-intro}.
We will also observe that, if $\CSPf{S} \notin \mNC^2$, then $\CSPf{S}$ is
$\pL$-hard.

\begin{theorem}
    \label{thm:cspsat-hard}
    Let $S$ be a finite set of Boolean relations.
    \begin{enumerate}
        \item 
            If
            $\CSPf{S} \notin 
            \mAC^0_3$
            then $\CSPf{S}$ is $\L$-hard
            under
            $\acmred$ reductions.
        \item 
            If
            $\CSPf{S} \notin \mNC^2$,
            then $\CSPf{S}$ is $\pL$-hard
            under
            $\acmred$ reductions.
    \end{enumerate}
\end{theorem}
\begin{proof}
    Item (1) follows 
            from \Cref{thm:ac0-csp}.
    To prove item (2),
            suppose that $\mdepth{\CSPf{S}} = {\omega(\log^2 n)}$.
            Then, by \Cref{thm:monfml-dichotomy},
            we conclude that 
            $\Pol(S) \sseq \clonefont{L_3}$, or
            $\Pol(S) \sseq \clonefont{V_2}$, or
            $\Pol(S) \sseq \clonefont{E_2}$.
            \Cref{thm:refined-completeness} 
            implies that
            $\CSPf{S}$ is $\pL$-hard.
            \qedhere
\end{proof}

\noindent \textbf{Further Discussion.} We recall the discussion of \Cref{sec:intro_mon_csp}.
We introduced and defined the functions $\CSPf{S}$ in that section,
as a way to capture monotone circuit lower bounds proved via lifting.
This in particular captures the monotone function $\txorsat$,
which was proved in~\cite{gkrs19} to require monotone circuit lower bounds
of size $2^{n^{\Omega(1)}}$ to compute, even though 
$\pL$-machines running in 
polynomial-time
can compute it.
\Cref{thm:cspsat-hard}
proves that this separation between monotone and non-monotone circuit lower
bounds cannot be improved by varying the set of relations $S$, as we argue
below.

There are two ways one could try to find a function in $\AC^0$
with large monotone complexity using a $\CSPSAT$ function.
First, one could try to define a set of relations $S$
such that $\CSPf{S} \in \AC^0$, but the monotone complexity of
$\CSPf{S}$ is large.
However, 
Item (1) of \Cref{thm:cspsat-hard}
proves that this is impossible,
as any $\CSPSAT$ function outside of $\mAC^0$
is $\L$-hard under simple reductions and, therefore, cannot be computed in $\AC^0$.

Secondly, one could try to be apply the arguments of
\Cref{s:constant-depth},
consisting of a padding trick and a simulation
theorem.
When $S$ is the set of 3XOR relations,
then indeed we obtain a function in $\AC^0[\xor]$ with superpolynomial
monotone circuit complexity,
as proved in \Cref{thm:xor-sat}.
However, 
Item (2) of \Cref{thm:cspsat-hard}
proves that this is best possible,
as any $\CSPSAT$ function which admits a superpolynomial monotone circuit
lower bound must be $\pL$-hard and, therefore,
at least as hard as $\txorsat$ for non-monotone circuits.
Item (2) also shows that even $\CSPSAT$ functions with a $\omega(\log^2 n)$
monotone depth lower bound must be $\pL$-hard,
which suggests that the arguments of
\Cref{s:constant-depth} applied to a $\CSPSAT$ function
are not able to prove the separation of \Cref{thm:ac0-fml-graph}.

A caveat to these impossibility results is in order.
First, 
we only study Boolean-valued CSPs here,
though the framework of lifting can also be applied in the context of
non-Boolean CSPs.
It's not clear if non-Boolean CSPs exhibit the same dichotomies
for monotone computation
we proved in this section.
We remark that Schaefer's dichotomy for Boolean-valued CSPs~\cite{DBLP:conf/stoc/Schaefer78}
has been extended to non-Boolean
CSPs~\cite{DBLP:conf/focs/Zhuk17,DBLP:conf/focs/Bulatov17}.

Secondly,
the instances of $\CSPSAT$ generated by lifting
do not cover the entirety of the minterms and maxterms of $\CSPSAT$.
In particular, our results do not rule out the possibility
that a clever interpolation of the instances generated by lifting
may give rise to a function that is easier to compute by non-monotone
circuits,
and therefore bypasses the hardness results of~\Cref{thm:cspsat-hard}.
One example is the Tardós function~\cite{tardos_1988}.
A lifting theorem applied to a Pigeonhole Principle formula can be used to prove 
a lower bound on the size of monotone circuits that accept cliques of size
$k$ and reject graphs that are $(k-1)$-colorable, for some $k = n^\eps$~\cite{DBLP:journals/combinatorica/RazM99,susanna_2023}.
A natural interpolation for these instances would be the $\kclq$
function, which, being $\NP$-complete, would be related to an $\NP$-complete
$\CSPSAT$.
However, as proved by~\cite{tardos_1988}, there is a monotone function
in
$\P$ which has the same output behaviour over these instances. 










\section{A Lower Bound for $\txorsat$ Using the Approximation Method}
\label{sec:xorsat-appr}

As discussed in \Cref{sec:intro_mon_csp},
\citep{gkrs19}
obtained an exponential lower bound on the monotone
circuit size of the function $\txorsat$ using techniques from communication
complexity and lifting. Here we observe that a weaker but still
super-polynomial lower bound  can be proved using the approximation method.

First, we recall the
function $\OddFactor_n : \blt^{\binom{n}{2}} \to \blt$
of \Cref{s:padding_graph},
which accepts 
a 
given graph
if the graph contains an \emph{odd factor},
which is a spanning subgraph in which the degree of every
vertex is odd.
For convenience,
in this section we consider a weaker version of $\OddFactor$, which takes as an input a
\emph{bipartite} graph with $n$ vertices on each part,
and accepts if the graph contains an odd factor.
Let $\BipOddFactor_n : \blt^{n^2} \to \blt$
be this function.
We remark that the lower bounds of
Babai, Gál and Wigderson~\cite{DBLP:journals/combinatorica/BabaiGW99}
for $\OddFactor$ (\Cref{thm:bip-oddfactor-lb})
also hold for $\BipOddFactor$.
The proof of the monotone circuit lower bound in particular
is essentially Razborov's lower bound for $\Matching$ via the approximation method~\cite{razborov1985lower}.

\begin{theorem}[\cite{DBLP:journals/combinatorica/BabaiGW99}]
    \label{thm:oddfactor}
    We have
    $$\Cmon{\BipOddFactor_n} = n^{\Omega(\log n)}
    \quad \text{and} \quad 
    \mdepth{\BipOddFactor_n} = \Omega(n).$$
\end{theorem}


We can reduce 
$\BipOddFactor$
to $\txorsat$ by
noting that computing $\BipOddFactor_n(M)$ on a given matrix
$M \in \blt^{n^2}$ is computationally equivalent to deciding the satisfiability of the following $\mathbb{F}_2$ linear system over variables $\{x_{ij}\}$:
\begin{itemize}
    \item For all $i \in [n]$:
        $\bigxor_{k=1}^n x_{ik} = 1$;
    \item For all $j \in [n]$:
        $\bigxor_{k=1}^n x_{kj} = 1$;
    \item For all $i,j \in [n]$ such that $M_{ij} = 0$:
        $x_{ij} = 0$.
\end{itemize}

We can then use a circuit for $\txorsat$ to solve this system
by using a standard trick of introducing new variables
to reduce the number of variables that appear in each equation,
as done in the textbook reduction from $\SAT$ to 3-$\SAT$. As the corresponding reductions turn out to be monotone, this implies monotone circuit and formula lower bounds for $\txorsat$.
We note that a somewhat similar   argument (in the non-monotone setting) appears in Feder and
Vardi~\cite[Theorem 30]{DBLP:journals/siamcomp/FederV98} regarding
constraint satisfaction problems with the ability to count.

In order to formalise this argument, we will need the following definition and results.

\begin{definition}
    \label{def:dual}
    Let $f$ be a Boolean function.
    We define $\dual(f) : x \mapsto \neg f(\neg x)$
    as the \emph{dual} of $f$.
\end{definition}

\begin{lemma}
    \label{lem:dual}
    Let $f$ be a monotone Boolean function.
    We have
    $\Cmon{f} = \Cmon{\dual(f)}$
    and
    $\mdepth{f} = \mdepth{\dual(f)}$.
\end{lemma}
\begin{proof}
    The idea is to push negations to the bottom and eliminate double negations at the input layer. In other words, applying De Morgan rules, we can convert any $\set{\land,\lor}$-circuit computing
    $f$ into a circuit computing $\dual(f)$ 
    by swapping $\land$-gates for
    $\lor$-gates, and vice-versa.
    Moreover, this transformation preserves the depth of the circuit.
\end{proof}




We are ready to describe a monotone reduction from the function $\BipOddFactor_n$ 
to $\txorsat$, which implies the desired lower bounds.

\begin{theorem}
    \label{thm:xorsat-lb}
    There exists $\eps > 0$
    such that
    $$\Cmon{\txorsat} = n^{\Omega(\log n)}
    \quad \text{and} \quad 
    \mdepth{\txorsat} = \Omega(n^\eps).$$
\end{theorem}
\begin{proof}
    Recall that the
    value of the function
    $\BipOddFactor_n(M)$ on a given matrix
    $M \in \blt^{n^2}$ is equal to 1 if the following system is
    satisfiable:
    \begin{itemize}
        \item For all $i \in [n]$:
            $\bigxor_{k=1}^n x_{ik} = 1$;
        \item For all $j \in [n]$:
            $\bigxor_{k=1}^n x_{kj} = 1$;
        \item For all $i,j \in [n]$ such that $M_{ij} = 0$:
            $x_{ij} = 0$.
    \end{itemize}
    We introduce some extra variables to reduce the number of variables in
    each equation in the following way.
    For every $i \in [n]$,
    introduce variables 
    $z_{i1},\dots,z_{i(n-1)}$
    and the equations
    \begin{align*}
        z_{i1} 
        &= 
        x_{i1} \xor x_{i2},
        \\
        z_{i2} 
        &= 
        z_{i1} \xor x_{i3},
        \\
        \dots
        \\
        z_{i,(n-1)} 
        &= 
        z_{i,(n-2)} \xor x_{i,n},
        \\
        z_{i,(n-1)}
        &= 
        1.
    \end{align*}
    Now note that these equations imply
    $z_{i,(n-2)} = \bigxor_{k=1}^n x_{ik} = 1$.
    For each ``column'' equation
    $\bigxor_{k=1}^n x_{kj} = 1$,
    we also add variables
    $w_{j1},\dots, w_{j(n-1)}$ as above.
    In total, we add
    at most
    $2n^2$ variables
    and
    $2n^2$ equations.
    Therefore, there is a system of linear equations
    on $O(n^2)$ variables, where each constraint contains at most $3$ variables, which is satisfiable
    if and only if
    $\BipOddFactor_n(M)=1$.
    Moreover, it is easy to see that the characteristic vector $\alpha$
    of the set of 
    equations of this system
    can be computed from $M$
    by an \emph{anti-monotone} projection, as we activate a constraint that depends on the input when $M_{ij} = 0$.

    Now let 
    $f = \dual(\txorsat)$ and $\beta = \neg \alpha$. Since, by definition, $\txorsat$ accepts \emph{unsatisfiable} systems,
    we get $\BipOddFactor_n(M) = \neg \txorsat(\alpha) = f(\beta)$ and
    that $\beta$ is a \emph{monotone} projection of $M$.
    Therefore, by \cref{lem:dual},
    we obtain
    \begin{equation*}
        \Cmon{\BipOddFactor_n} \leq \Cmon{\txorsat}
    \end{equation*}
    and
    \begin{equation*}
        \mdepth{\BipOddFactor_n} \leq \mdepth{\txorsat}.
        \qedhere
    \end{equation*}
\end{proof}


\section{Schaefer's Theorem in Monotone Complexity}
\label{sec:schaefer_proof}

\subsection{Connectivity and generation functions}

We recall the definitions of two prominent monotone Boolean functions
that have efficient monotone circuits.
Let $\stconn : \blt^{n^2} \to \blt$
be the function that outputs $1$ on a given directed graph $G$
if there exists a path from $1$ to $n$ in $G$.
Let $\GEN : \blt^{n^3} \to \blt$
be the Boolean function which receives a set $T$ of triples 
$(i,j,k) \in [n^3]$,
and outputs $1$ if 
$n \in S$, where $S \sseq [n]$ is the 
set generated with the following rules:

\begin{itemize}
    \item \emph{Axiom:} $1 \in S$,
    \item \emph{Generation:} If $i,j \in S$ and $(i,j,k) \in T$, then
        $k \in S$.
\end{itemize}

The following upper bounds are well-known and easy to prove.

\begin{theorem}
    [\protect{\cite[Exercise 7.3]{jukna_2012}, \cite{DBLP:journals/combinatorica/RazM99}}]
    \label{thm:stconn-gen}
    We have $\stconn \in \mNL$ and
    $\GEN \in \mSIZE[\poly]$.
\end{theorem}

\subsection{Proof of reduction lemmas}
\label{sec:reduction_monotone}

Here we present monotonised versions of the proofs of
\cite[Propositions 2.2 - 2.4]{playing_two},
which give a simplified presentation of the results
of~\cite{DBLP:conf/stoc/Schaefer78}.


\inclcoq*

\begin{proof}
    If $\COQ(S_1) \sseq \COQ(S_2)$, then
    each relation of $S_1$ can be represented as a conjunctive query
    over $S_2$.
    Let $F_1$ be a $S_1$-formula.
    For each constraint $C_1$ of $F_1$, there exists
    a formula $\varphi(C_1)$ in $\CNF(S_2)$
    such that
    $C_1$ is a projection of $\varphi(C_1)
    $ (i.e., $C_1$ is a conjunctive query of
    $\varphi(C_1)$).
    However, note that $C_1$ is satisfiable if and only if
    $\varphi(C_1)$ is satisfiable.
    So we can replace the constraint $C_1$ by the \emph{set} of constraints
    in $\varphi(C_1)$.
    Doing this for every constraint in $F_1$, we obtain an $S_2$-formula $F_2$
    which is satisfiable iff $F_1$ is satisfiable.

    Now note that, to decide if a given constraint application $C$ of
    $S_2$
    is in the reduction,
    it suffices
    to check if there exists a $S_1$-constraint $C_1$
    in $F_1$
    such that
    $C$ is in $\varphi(C_1)$.
    Using non-uniformity, this can be easily done
    by an OR over the relevant input bits.

    Finally, we observe that, since the arities of each relation in
    $S_1$ and $S_2$ are constant, we only add a constant number of
    variables for each constraint
    to represent $S_1$-formulas with conjunctive queries over $S_2$-formulas.
\end{proof}


\begin{lemma}
    \label{prop:eql}
    Let $S$ be a set of Boolean relations.
    We have
    $\CSPf{S \cup \set{=}} \mnlmred \CSPf{S}$.
\end{lemma}
\begin{proof}
    Let $F$ be a $(S \cup \set{=})$-formula on $n$ variables
    given as an input.
    Remember that $F$ is given as a Boolean vector $\alpha$,
    where each bit of $\alpha$
    represents
    the presence of a constraint application on $n$ variables
    from $S \cup \set{=}$.
    We first build an undirected graph $G$ with the variables
    $x_1,\dots,x_n$ as vertices,
    and we put an edge between $x_i$ and $x_j$
    if the constraint $x_i = x_j$ appears in $F$.
    Note that $G$ can be constructed by a monotone projection from
    $F$.

    Let $R \in S$ and let
    $C = R(x_1,\dots,x_n)$ be a constraint application of $R$.
    If $C$ appears in $F$,
    we add to the system every constraint of the form
    $C' = R(y_1,\dots,y_n)$
    such that,
    for every $i \in [n]$,
    there exists a path from $x_i$ to $y_i$ in the graph $G$.
    In this case, we say that $C$ \emph{generates} $C'$.
    Let $F_2$ be the formula that contains all
    non-equality constraints of $F$,
    and all the non-equality constraints
    generated by a constraint in $F$.
    It's not hard to see that 
    $F$ is satisfiable if and only if $F_2$ is satisfiable,
    and therefore the reduction is correct.

    Moreover, the reduction can be done in monotone $\NL$
    using the monotone $\NL$ algorithm for $\stconn$
    (\Cref{thm:stconn-gen}).
    Indeed, there are at most $n^k$ constraint applications of a given
    relation $R$ of arity $k$. Therefore,
    to decide if a constraint $C' = R(y_1,\dots,y_n)$ appears in $F_2$,
    it suffices to check if there exists a constraint application 
    of $R$ in $F$ which generates $C'$.
    This can be checked with $n^k$ calls to $\stconn$.
\end{proof}

\polym*
\begin{proof}
    If $\Pol(S_2) \sseq \Pol(S_1)$, then 
    from \Cref{lem:galois}
    (\Cref{item:inv-invert,item:coq-sseq,item:invpol})
    we obtain
    $\COQ(S_1) \sseq \ccln{S_1} \sseq \ccln{S_2} = \COQ(S_2 \cup \set{=})$.
    Therefore, by~\Cref{prop:incl-coq,prop:eql} 
    we can do the following chain of reductions in monotone $\NL$:
    \begin{equation*}
        \CSPf{S_1}
        \mormred
        \CSPf{S_2 \cup \set{=}}
        \mnlmred
        \CSPf{S_2}.
        \qedhere
    \end{equation*}
\end{proof}


\subsection{Monotone circuit upper bounds}
\label{sec:monotone_upper}





We restate and prove the theorem.

\monschaeffer*

\begin{proof}
    We prove each case separately.

    \emph{Proof of (1)}.
    We first observe that
    $\thornsat$ 
    (see definition in \Cref{sec:csp-dichotomy}, \emph{Dichotomy for formulas})
    can be solved by a reduction to
    $\GEN \in \mSIZE[\poly]$.
    Indeed, we interpret each constraint of the form
    $(\neg x_i \lor \neg x_j \lor x_k)$
    (which is equivalent to 
    $x_i \land x_j \implies x_k$)
    as a triple
    $(i,j,k)$,
    and constraints of the form
    $x_i$
    as a triple
    $(0,0,i)$.
    Let $S \sseq \set{0,1,2\dots,n}$ be the set generated
    by these triples, applying the generation rules of $\GEN$,
    using $0 \in S$ as the axiom.
    It suffices to check that there exists
    some constraint of the form
    $\neg x_i \lor \neg x_j \lor \neg x_k$,
    such that $\set{i,j,k} \sseq S$.
    This process can be done with polynomial-size
    monotone circuits, invoking $\GEN$.
    Therefore, it follows
    from \Cref{thm:stconn-gen}
    that
    $\thornsat \in \mSIZE[\poly]$.

    Moreover, we recall that, if $\clonefont{E_2} \sseq \Pol(S)$, then
    $S \sseq \COQ(\hornt)$
    (in other words, every $S$-formula can be written as a set of 3-Horn
    equations) -- see, e.g,~\cite[Lemma 4.8]{DBLP:books/daglib/0004131}.
    Therefore, from 
    \Cref{item:coq-idempotent,item:coq-sseq} of \Cref{lem:galois}
    and
    \Cref{prop:incl-coq},
    we conclude that
    $\CSPf{S} \mormred \thornsat \in \mSIZE[\poly]$.

    Now recall that, if $\clonefont{V_2} \sseq \Pol(S)$,
    then
    $S \sseq \COQ(\ahornt)$,
    where
    $\ahornt$ is the set of width-3 Anti-Horn relations
    (i.e., $\ahornt = 
    \set{
        (x_1 \lor x_2 \lor \neg x_3),
        (x_1 \lor x_2 \lor x_3),
        (\neg x)
    }$; see~\cite[Lemma 4.8]{DBLP:books/daglib/0004131} for a proof of this
    observation).
    But note that an
    $\ahornt$-formula $\varphi$
    is satisfiable if and only if
    the $\hornt$-formula 
    $\varphi(\neg x_1,\dots,\neg x_n)$
    is satisfiable.
    Therefore by~\Cref{prop:incl-coq}
    and
    \Cref{item:coq-idempotent,item:coq-sseq} of \Cref{lem:galois},
    we have
    $\CSPf{S} \mormred 
    \CSPf{\ahornt}
    \mprojred \thornsat \in \mSIZE[\poly]$.

    \emph{Proof of (2)}.
    We first prove the case $\clonefont{D_2} \sseq \Pol(S)$.
    Let $\dSAT = \CSPf{\Gamma}$,
    where \[\Gamma = \set{(x_1 \lor x_2), (x_1 \lor \neg x_2), (\neg x_1 \lor
    \neg x_2)}\].
    It's easy to check that the standard reduction from
    $\dSAT$ to $\stconn$ can be done in monotone $\NL$
    (see~\cite[Theorem 4]{DBLP:journals/mst/JonesLL76}).
    Therefore, it follows from \Cref{thm:stconn-gen}
    that
    $\dSAT \in \mNL$.
    Now, recall that, if
    $\clonefont{D_2} \sseq \Pol(S)$, then
    $S \sseq \COQ(\Gamma)$
    (see, e.g., \cite[Lemma 4.9]{DBLP:books/daglib/0004131}).
    Therefore,
    from \Cref{prop:incl-coq} and
    \Cref{item:coq-idempotent,item:coq-sseq} of \Cref{lem:galois},
    we conclude
    $\CSPf{S} \in \mNL$.

    We now suppose that $\clonefont{S_{00}} \sseq \Pol(S)$.
    We check that the proof of~\cite[Lemma 3.4]{DBLP:journals/jcss/AllenderBISV09}
    gives a monotone circuit.
    If $\clonefont{S_{00}} \sseq \Pol(S)$, then
    there exists $k \geq 2$ such that
    $\clonefont{S_{00}}^k \sseq \Pol(S)$
    (that's because there does not exist a finite set of relations
    $S$ such that $\Pol(S) = \clonefont{S_{00}}$).
    Note that
    $\clonefont{S_{00}}^k = \Pol(\Gamma)$,
    where $\Gamma = \set{\OR^k, x, \neg x, \to, =}$.
    We show below how to decide if a $\Gamma$-formula is unsatisfiable 
    in monotone $\NL$.
    The result then follows from~\Cref{prop:poly}.

    Let $F$ be a given $\Gamma$-formula with $n$ variables.
    We first construct a directed graph $G$,
    with vertex set $\set{x_1,\dots,x_n}$,
    and with arcs 
    $(x_i,x_j)$
    if $x_i \to x_j$ is a constraint of $F$,
    and 
    arcs
    $(x_i,x_j)$
    and
    $(x_j,x_i)$
    if
    $x_i = x_j$ is a constraint of $F$.
    This can be done with a monotone projection.
    Observe that a $\Gamma$-formula $F$
    is unsatisfiable if, and only if,
    there exists a constraint of the form $x_{i_1} \lor \dots \lor x_{i_k}$ in $F$,
    such that there exists a path from some $x_{i_j}$ to a constraint
    $\neg y$ in $F$.
    This can be checked in monotone $\NL$ by \Cref{thm:stconn-gen}.

    The case $\clonefont{S_{10}} \sseq \Pol(S)$ is analogous.
\end{proof}



\section{Background on Post's Lattice and Clones}
\label{sec:clone-background}

In this section, we include
 the definitions of the various clones that are used in the
paper,
as well as
a figure of Post's lattice, which can be helpful when checking the proofs
of
\Cref{sec:CSPs}.

Let
$\rightarrow : (x,y) \mapsto (\neg x \lor y)$.
Let also
$\leftrightarrow : (x,y) \mapsto \neg (x \xor y)$
and
$\clonefont{id} : x \mapsto x$.
Let $f : \blt^k \to \blt$ be a Boolean function.
We say that 
$f$
is \emph{linear}
if there exists 
$c \in \blt^k$ and $b \in \blt$
such that
$f(x) = \inner{c}{x} + b \pmod 2$.
We say that $f$ is self-dual if $f = \dual(f)$. Let $a \in \blt$.
We say that $f$ is 
\emph{a-reproducing}
if
$f(a,\dots,a) = a$.
We say that
a set $T \sseq \blt^k$ is 
\emph{$a$-separating}
if there exists $i \in [k]$
such that
$x_i = a$ for all $x \in T$.
We say that
$f$ is 
\emph{$a$-separating}
if
$f^{-1}(a)$ is $a$-separating.
We say that $f$ is 
\emph{a-separating of degree $k$}
if every $T \sseq f^{-1}(a)$
such that $\card{T} = k$ is $a$-separating. The \emph{basis} of a clone $B$ is a set of Boolean functions $F$
such that $B = [F]$.

 


\begin{figure}[hbtp]
    \centering
    \begin{small}
        \begin{longtable}[]{p{0.1\textwidth}p{0.5\textwidth}p{0.3\textwidth}}
            \toprule
            \begin{minipage}[b]{0.12\columnwidth}\raggedright
                Name\strut
            \end{minipage} & \begin{minipage}[b]{0.41\columnwidth}\raggedright
                Definition\strut
            \end{minipage} & \begin{minipage}[b]{0.38\columnwidth}\raggedright
                Base\strut
            \end{minipage}\tabularnewline
            \midrule
            \endhead
            \begin{minipage}[t]{0.12\columnwidth}\raggedright
                \(\mathrm{BF}\)\strut
            \end{minipage} & \begin{minipage}[t]{0.41\columnwidth}\raggedright
                All Boolean functions\strut
            \end{minipage} & \begin{minipage}[t]{0.38\columnwidth}\raggedright
                \(\set{\lor,\land,\neg}\)\strut
            \end{minipage}\tabularnewline
            \begin{minipage}[t]{0.12\columnwidth}\raggedright
                \(\clonefont{R_0}\)\strut
            \end{minipage} & \begin{minipage}[t]{0.41\columnwidth}\raggedright
                \(\set{f \in \clonefont{BF} : f \text{ is $0$-reproducing}}\)\strut
            \end{minipage} & \begin{minipage}[t]{0.38\columnwidth}\raggedright
                \(\set{\land,\xor}\)\strut
            \end{minipage}\tabularnewline
            \begin{minipage}[t]{0.12\columnwidth}\raggedright
                \(\clonefont{R_1}\)\strut
            \end{minipage} & \begin{minipage}[t]{0.41\columnwidth}\raggedright
                \(\set{f \in \clonefont{BF} : f \text{ is $1$-reproducing}}\)\strut
            \end{minipage} & \begin{minipage}[t]{0.38\columnwidth}\raggedright
                \(\set{\lor,\leftrightarrow}\)\strut
            \end{minipage}\tabularnewline
            \begin{minipage}[t]{0.12\columnwidth}\raggedright
                \(\clonefont{R_2}\)\strut
            \end{minipage} & \begin{minipage}[t]{0.41\columnwidth}\raggedright
                \(\clonefont{R_1 \cap R_0}\)\strut
            \end{minipage} & \begin{minipage}[t]{0.38\columnwidth}\raggedright
                \(\set{\lor, x \land (y \leftrightarrow z)}\)\strut
            \end{minipage}\tabularnewline
            \begin{minipage}[t]{0.12\columnwidth}\raggedright
                \(\clonefont{M}\)\strut
            \end{minipage} & \begin{minipage}[t]{0.41\columnwidth}\raggedright
                \(\set{f \in \clonefont{BF} : f \text{ is monotonic}}\)\strut
            \end{minipage} & \begin{minipage}[t]{0.38\columnwidth}\raggedright
                \(\set{\lor,\land,0,1}\)\strut
            \end{minipage}\tabularnewline
            \begin{minipage}[t]{0.12\columnwidth}\raggedright
                \(\clonefont{M_1}\)\strut
            \end{minipage} & \begin{minipage}[t]{0.41\columnwidth}\raggedright
                \(\clonefont{M} \cap \clonefont{R_1}\)\strut
            \end{minipage} & \begin{minipage}[t]{0.38\columnwidth}\raggedright
                \(\set{\lor,\land,1}\)\strut
            \end{minipage}\tabularnewline
            \begin{minipage}[t]{0.12\columnwidth}\raggedright
                \(\clonefont{M_0}\)\strut
            \end{minipage} & \begin{minipage}[t]{0.41\columnwidth}\raggedright
                \(\clonefont{M} \cap \clonefont{R_0}\)\strut
            \end{minipage} & \begin{minipage}[t]{0.38\columnwidth}\raggedright
                \(\set{\lor,\land,0}\)\strut
            \end{minipage}\tabularnewline
            \begin{minipage}[t]{0.12\columnwidth}\raggedright
                \(\clonefont{M_2}\)\strut
            \end{minipage} & \begin{minipage}[t]{0.41\columnwidth}\raggedright
                \(\clonefont{M} \cap \clonefont{R_2}\)\strut
            \end{minipage} & \begin{minipage}[t]{0.38\columnwidth}\raggedright
                \(\set{\lor,\land}\)\strut
            \end{minipage}\tabularnewline
            \begin{minipage}[t]{0.12\columnwidth}\raggedright
                \(\clonefont{S_0^n}\)\strut
            \end{minipage} & \begin{minipage}[t]{0.41\columnwidth}\raggedright
                \(\set{f \in \clonefont{BF} : f \text{ is $0$-separating of degree $n$}}\)\strut
            \end{minipage} & \begin{minipage}[t]{0.38\columnwidth}\raggedright
                \(\set{\rightarrow,\dual(h_n)}\)\strut
            \end{minipage}\tabularnewline
            \begin{minipage}[t]{0.12\columnwidth}\raggedright
                \(\clonefont{S_0}\)\strut
            \end{minipage} & \begin{minipage}[t]{0.41\columnwidth}\raggedright
                \(\set{f \in \clonefont{BF} : f \text{ is $0$-separating}}\)\strut
            \end{minipage} & \begin{minipage}[t]{0.38\columnwidth}\raggedright
                \(\set{\rightarrow}\)\strut
            \end{minipage}\tabularnewline
            \begin{minipage}[t]{0.12\columnwidth}\raggedright
                \(\clonefont{S_1^n}\)\strut
            \end{minipage} & \begin{minipage}[t]{0.41\columnwidth}\raggedright
                \(\set{f \in \clonefont{BF} : f \text{ is $1$-separating of degree $n$}}\)\strut
            \end{minipage} & \begin{minipage}[t]{0.38\columnwidth}\raggedright
                \(\set{x \land \ovl{y} ,h_n}\)\strut
            \end{minipage}\tabularnewline
            \begin{minipage}[t]{0.12\columnwidth}\raggedright
                \(\clonefont{S_1}\)\strut
            \end{minipage} & \begin{minipage}[t]{0.41\columnwidth}\raggedright
                \(\set{f \in \clonefont{BF} : f \text{ is $1$-separating}}\)\strut
            \end{minipage} & \begin{minipage}[t]{0.38\columnwidth}\raggedright
                \(\set{x \land \ovl{y}}\)\strut
            \end{minipage}\tabularnewline
            \begin{minipage}[t]{0.12\columnwidth}\raggedright
                \(\clonefont{S_{02}^n}\)\strut
            \end{minipage} & \begin{minipage}[t]{0.41\columnwidth}\raggedright
                \(\clonefont{S_0^n} \cap \clonefont{R_2}\)\strut
            \end{minipage} & \begin{minipage}[t]{0.38\columnwidth}\raggedright
                \(\set{x \lor (y \land \ovl{z}), \dual(h_n)}\)\strut
            \end{minipage}\tabularnewline
            \begin{minipage}[t]{0.12\columnwidth}\raggedright
                \(\clonefont{S_{02}}\)\strut
            \end{minipage} & \begin{minipage}[t]{0.41\columnwidth}\raggedright
                \(\clonefont{S_0} \cap \clonefont{R_2}\)\strut
            \end{minipage} & \begin{minipage}[t]{0.38\columnwidth}\raggedright
                \(\set{x \lor (y \land \ovl{z})}\)\strut
            \end{minipage}\tabularnewline
            \begin{minipage}[t]{0.12\columnwidth}\raggedright
                \(\clonefont{S_{01}^n}\)\strut
            \end{minipage} & \begin{minipage}[t]{0.41\columnwidth}\raggedright
                \(\clonefont{S_0^n} \cap \clonefont{M}\)\strut
            \end{minipage} & \begin{minipage}[t]{0.38\columnwidth}\raggedright
                \(\set{\dual(h_n), 1}\)\strut
            \end{minipage}\tabularnewline
            \begin{minipage}[t]{0.12\columnwidth}\raggedright
                \(\clonefont{S_{01}}\)\strut
            \end{minipage} & \begin{minipage}[t]{0.41\columnwidth}\raggedright
                \(\clonefont{S_0} \cap \clonefont{M}\)\strut
            \end{minipage} & \begin{minipage}[t]{0.38\columnwidth}\raggedright
                \(\set{x \lor (y \land z), 1}\)\strut
            \end{minipage}\tabularnewline
            \begin{minipage}[t]{0.12\columnwidth}\raggedright
                \(\clonefont{S_{00}^n}\)\strut
            \end{minipage} & \begin{minipage}[t]{0.41\columnwidth}\raggedright
                \(\clonefont{S_0^n} \cap \clonefont{R_2} \cap \clonefont{M}\)\strut
            \end{minipage} & \begin{minipage}[t]{0.38\columnwidth}\raggedright
                \(\set{x \lor (y \land z), \dual(h_n)}\)\strut
            \end{minipage}\tabularnewline
            \begin{minipage}[t]{0.12\columnwidth}\raggedright
                \(\clonefont{S_{00}}\)\strut
            \end{minipage} & \begin{minipage}[t]{0.41\columnwidth}\raggedright
                \(\clonefont{S_0} \cap \clonefont{R_2} \cap \clonefont{M}\)\strut
            \end{minipage} & \begin{minipage}[t]{0.38\columnwidth}\raggedright
                \(\set{x \lor (y \land z)}\)\strut
            \end{minipage}\tabularnewline
            \begin{minipage}[t]{0.12\columnwidth}\raggedright
                \(\clonefont{S_{12}^n}\)\strut
            \end{minipage} & \begin{minipage}[t]{0.41\columnwidth}\raggedright
                \(\clonefont{S_1^n} \cap \clonefont{R_2}\)\strut
            \end{minipage} & \begin{minipage}[t]{0.38\columnwidth}\raggedright
                \(\set{x \land (y \lor \ovl{z}), h_n}\)\strut
            \end{minipage}\tabularnewline
            \begin{minipage}[t]{0.12\columnwidth}\raggedright
                \(\clonefont{S_{12}}\)\strut
            \end{minipage} & \begin{minipage}[t]{0.41\columnwidth}\raggedright
                \(\clonefont{S_1} \cap \clonefont{R_2}\)\strut
            \end{minipage} & \begin{minipage}[t]{0.38\columnwidth}\raggedright
                \(\set{x \land (y \lor \ovl{z})}\)\strut
            \end{minipage}\tabularnewline
            \begin{minipage}[t]{0.12\columnwidth}\raggedright
                \(\clonefont{S_{11}^n}\)\strut
            \end{minipage} & \begin{minipage}[t]{0.41\columnwidth}\raggedright
                \(\clonefont{S_1^n} \cap \clonefont{M}\)\strut
            \end{minipage} & \begin{minipage}[t]{0.38\columnwidth}\raggedright
                \(\set{h_n, 0}\)\strut
            \end{minipage}\tabularnewline
            \begin{minipage}[t]{0.12\columnwidth}\raggedright
                \(\clonefont{S_{11}}\)\strut
            \end{minipage} & \begin{minipage}[t]{0.41\columnwidth}\raggedright
                \(\clonefont{S_1} \cap \clonefont{M}\)\strut
            \end{minipage} & \begin{minipage}[t]{0.38\columnwidth}\raggedright
                \(\set{x \land (y \lor z), 0}\)\strut
            \end{minipage}\tabularnewline
            \begin{minipage}[t]{0.12\columnwidth}\raggedright
                \(\clonefont{S_{10}^n}\)\strut
            \end{minipage} & \begin{minipage}[t]{0.41\columnwidth}\raggedright
                \(\clonefont{S_1^n} \cap \clonefont{R_2} \cap \clonefont{M}\)\strut
            \end{minipage} & \begin{minipage}[t]{0.38\columnwidth}\raggedright
                \(\set{x \land (y \lor z), h_n}\)\strut
            \end{minipage}\tabularnewline
            \begin{minipage}[t]{0.12\columnwidth}\raggedright
                \(\clonefont{S_{10}}\)\strut
            \end{minipage} & \begin{minipage}[t]{0.41\columnwidth}\raggedright
                \(\clonefont{S_1} \cap \clonefont{R_2} \cap \clonefont{M}\)\strut
            \end{minipage} & \begin{minipage}[t]{0.38\columnwidth}\raggedright
                \(\set{x \land (y \lor z)}\)\strut
            \end{minipage}\tabularnewline
            \begin{minipage}[t]{0.12\columnwidth}\raggedright
                \(\clonefont{D}\)\strut
            \end{minipage} & \begin{minipage}[t]{0.41\columnwidth}\raggedright
                \(\set{f \in \clonefont{BF} : f \text{ is self-dual}}\)\strut
            \end{minipage} & \begin{minipage}[t]{0.38\columnwidth}\raggedright
                \(\set{(x \land \ovl{y}) \lor (x \land \ovl{z}) \lor (\ovl{y} \land \ovl{z})}\)\strut
            \end{minipage}\tabularnewline
            \begin{minipage}[t]{0.12\columnwidth}\raggedright
                \(\clonefont{D_1}\)\strut
            \end{minipage} & \begin{minipage}[t]{0.41\columnwidth}\raggedright
                \(\clonefont{D} \cap \clonefont{R_2}\)\strut
            \end{minipage} & \begin{minipage}[t]{0.38\columnwidth}\raggedright
                \(\set{(x \land y) \lor (x \land \ovl{z}) \lor ({y} \land \ovl{z})}\)\strut
            \end{minipage}\tabularnewline
            \begin{minipage}[t]{0.12\columnwidth}\raggedright
                \(\clonefont{D_2}\)\strut
            \end{minipage} & \begin{minipage}[t]{0.41\columnwidth}\raggedright
                \(\clonefont{D} \cap \clonefont{M}\)\strut
            \end{minipage} & \begin{minipage}[t]{0.38\columnwidth}\raggedright
                \(\set{(x \land y) \lor (y \land {z}) \lor (x \land {z})}\)\strut
            \end{minipage}\tabularnewline
            \begin{minipage}[t]{0.12\columnwidth}\raggedright
                \(\clonefont{L}\)\strut
            \end{minipage} & \begin{minipage}[t]{0.41\columnwidth}\raggedright
                \(\set{f \in \clonefont{BF} : f \text{ is linear}}\)\strut
            \end{minipage} & \begin{minipage}[t]{0.38\columnwidth}\raggedright
                \(\set{\xor, 1}\)\strut
            \end{minipage}\tabularnewline
            \begin{minipage}[t]{0.12\columnwidth}\raggedright
                \(\clonefont{L_0}\)\strut
            \end{minipage} & \begin{minipage}[t]{0.41\columnwidth}\raggedright
                \(\clonefont{L} \cap \clonefont{R_0}\)\strut
            \end{minipage} & \begin{minipage}[t]{0.38\columnwidth}\raggedright
                \(\set{\xor}\)\strut
            \end{minipage}\tabularnewline
            \begin{minipage}[t]{0.12\columnwidth}\raggedright
                \(\clonefont{L_1}\)\strut
            \end{minipage} & \begin{minipage}[t]{0.41\columnwidth}\raggedright
                \(\clonefont{L} \cap \clonefont{R_1}\)\strut
            \end{minipage} & \begin{minipage}[t]{0.38\columnwidth}\raggedright
                \(\set{\leftrightarrow}\)\strut
            \end{minipage}\tabularnewline
            \begin{minipage}[t]{0.12\columnwidth}\raggedright
                \(\clonefont{L_2}\)\strut
            \end{minipage} & \begin{minipage}[t]{0.41\columnwidth}\raggedright
                \(\clonefont{L} \cap \clonefont{R}\)\strut
            \end{minipage} & \begin{minipage}[t]{0.38\columnwidth}\raggedright
                \(\set{x \xor y \xor z}\)\strut
            \end{minipage}\tabularnewline
            \begin{minipage}[t]{0.12\columnwidth}\raggedright
                \(\clonefont{L_3}\)\strut
            \end{minipage} & \begin{minipage}[t]{0.41\columnwidth}\raggedright
                \(\clonefont{L} \cap \clonefont{D}\)\strut
            \end{minipage} & \begin{minipage}[t]{0.38\columnwidth}\raggedright
                \(\set{x \xor y \xor z \xor 1}\)\strut
            \end{minipage}\tabularnewline
            \begin{minipage}[t]{0.12\columnwidth}\raggedright
                \(\clonefont{V}\)\strut
            \end{minipage} & \begin{minipage}[t]{0.41\columnwidth}\raggedright
                \(\set{f \in \clonefont{BF} : f \text{ is constant or an $n$-ary OR function}}\)\strut
            \end{minipage} & \begin{minipage}[t]{0.38\columnwidth}\raggedright
                \(\set{\lor,0,1}\)\strut
            \end{minipage}\tabularnewline
            \begin{minipage}[t]{0.12\columnwidth}\raggedright
                \(\clonefont{V_0}\)\strut
            \end{minipage} & \begin{minipage}[t]{0.41\columnwidth}\raggedright
                \([\set{\lor}] \cup [\set{0}]\)\strut
            \end{minipage} & \begin{minipage}[t]{0.38\columnwidth}\raggedright
                \(\set{\lor,0}\)\strut
            \end{minipage}\tabularnewline
            \begin{minipage}[t]{0.12\columnwidth}\raggedright
                \(\clonefont{V_1}\)\strut
            \end{minipage} & \begin{minipage}[t]{0.41\columnwidth}\raggedright
                \([\set{\lor}] \cup [\set{1}]\)\strut
            \end{minipage} & \begin{minipage}[t]{0.38\columnwidth}\raggedright
                \(\set{\lor,1}\)\strut
            \end{minipage}\tabularnewline
            \begin{minipage}[t]{0.12\columnwidth}\raggedright
                \(\clonefont{V_2}\)\strut
            \end{minipage} & \begin{minipage}[t]{0.41\columnwidth}\raggedright
                \([\set{\lor}]\)\strut
            \end{minipage} & \begin{minipage}[t]{0.38\columnwidth}\raggedright
                \(\set{\lor}\)\strut
            \end{minipage}\tabularnewline
            \begin{minipage}[t]{0.12\columnwidth}\raggedright
                \(\clonefont{E}\)\strut
            \end{minipage} & \begin{minipage}[t]{0.41\columnwidth}\raggedright
                \(\set{f \in \clonefont{BF} : f \text{ is constant or an $n$-ary AND function}}\)\strut
            \end{minipage} & \begin{minipage}[t]{0.38\columnwidth}\raggedright
                \(\set{\land,0,1}\)\strut
            \end{minipage}\tabularnewline
            \begin{minipage}[t]{0.12\columnwidth}\raggedright
                \(\clonefont{E_0}\)\strut
            \end{minipage} & \begin{minipage}[t]{0.41\columnwidth}\raggedright
                \([\set{\land}] \cup [\set{0}]\)\strut
            \end{minipage} & \begin{minipage}[t]{0.38\columnwidth}\raggedright
                \(\set{\land,0}\)\strut
            \end{minipage}\tabularnewline
            \begin{minipage}[t]{0.12\columnwidth}\raggedright
                \(\clonefont{E_1}\)\strut
            \end{minipage} & \begin{minipage}[t]{0.41\columnwidth}\raggedright
                \([\set{\land}] \cup [\set{1}]\)\strut
            \end{minipage} & \begin{minipage}[t]{0.38\columnwidth}\raggedright
                \(\set{\land,1}\)\strut
            \end{minipage}\tabularnewline
            \begin{minipage}[t]{0.12\columnwidth}\raggedright
                \(\clonefont{E_2}\)\strut
            \end{minipage} & \begin{minipage}[t]{0.41\columnwidth}\raggedright
                \([\set{\land}]\)\strut
            \end{minipage} & \begin{minipage}[t]{0.38\columnwidth}\raggedright
                \(\set{\land}\)\strut
            \end{minipage}\tabularnewline
            \begin{minipage}[t]{0.12\columnwidth}\raggedright
                \(\clonefont{N}\)\strut
            \end{minipage} & \begin{minipage}[t]{0.41\columnwidth}\raggedright
                \([\set{\neg}] \cup [\set{0}] \cup [\set{1}]\)\strut
            \end{minipage} & \begin{minipage}[t]{0.38\columnwidth}\raggedright
                \(\set{\neg, 1}\)\strut
            \end{minipage}\tabularnewline
            \begin{minipage}[t]{0.12\columnwidth}\raggedright
                \(\clonefont{N_2}\)\strut
            \end{minipage} & \begin{minipage}[t]{0.41\columnwidth}\raggedright
                \([\set{\neg}]\)\strut
            \end{minipage} & \begin{minipage}[t]{0.38\columnwidth}\raggedright
                \(\set{\neg}\)\strut
            \end{minipage}\tabularnewline
            \begin{minipage}[t]{0.12\columnwidth}\raggedright
                \(\clonefont{I}\)\strut
            \end{minipage} & \begin{minipage}[t]{0.41\columnwidth}\raggedright
                \([\set{\clonefont{id}}] \cup [\set{0}] \cup [\set{1}]\)\strut
            \end{minipage} & \begin{minipage}[t]{0.38\columnwidth}\raggedright
                \(\set{\clonefont{id}, 0, 1}\)\strut
            \end{minipage}\tabularnewline
            \begin{minipage}[t]{0.12\columnwidth}\raggedright
                \(\clonefont{I_0}\)\strut
            \end{minipage} & \begin{minipage}[t]{0.41\columnwidth}\raggedright
                \([\set{\clonefont{id}}] \cup [\set{0}]\)\strut
            \end{minipage} & \begin{minipage}[t]{0.38\columnwidth}\raggedright
                \(\set{\clonefont{id}, 0}\)\strut
            \end{minipage}\tabularnewline
            \begin{minipage}[t]{0.12\columnwidth}\raggedright
                \(\clonefont{I_1}\)\strut
            \end{minipage} & \begin{minipage}[t]{0.41\columnwidth}\raggedright
                \([\set{\clonefont{id}}] \cup [\set{1}]\)\strut
            \end{minipage} & \begin{minipage}[t]{0.38\columnwidth}\raggedright
                \(\set{\clonefont{id}, 1}\)\strut
            \end{minipage}\tabularnewline
            \begin{minipage}[t]{0.12\columnwidth}\raggedright
                \(\clonefont{I_2}\)\strut
            \end{minipage} & \begin{minipage}[t]{0.41\columnwidth}\raggedright
                \([\set{\clonefont{id}}]\)\strut
            \end{minipage} & \begin{minipage}[t]{0.38\columnwidth}\raggedright
                \(\set{\clonefont{id}}\)\strut
            \end{minipage}\tabularnewline
            \bottomrule
        \end{longtable}
    \end{small}
    \caption{
        Table of all closed classes of Boolean functions, and their bases.
        Here, $h_n$ denotes the function
        $h_n(x_1,\dots,x_{n+1}) = \bigvee_{i=1}^{n+1} \bigwedge_{j = 1,
        j \neq i}^{n+1} x_j$.
        See \Cref{def:dual} for the definition of $\dual(\cdot)$.
        The same
        table appears in
    \cite[Table 1]{DBLP:journals/jcss/AllenderBISV09}.}
    \label{fig:clone_table}
\end{figure}

\newcommand{\etalchar}[1]{$^{#1}$}
\begin{thebibliography}{HWWY94}

\bibitem[AB87]{DBLP:journals/combinatorica/AlonB87}
Noga Alon and Ravi~B. Boppana.
\newblock The monotone circuit complexity of boolean functions.
\newblock {\em Combinatorica}, 7(1):1--22, 1987.

\bibitem[ABI{\etalchar{+}}09]{DBLP:journals/jcss/AllenderBISV09}
Eric Allender, Michael Bauland, Neil Immerman, Henning Schnoor, and Heribert
  Vollmer.
\newblock The complexity of satisfiability problems: Refining {S}chaefer's
  theorem.
\newblock {\em J. Comput. Syst. Sci.}, 75(4):245--254, 2009.

\bibitem[AG87]{ajtai_gurevich_1987}
Mikl\'{o}s Ajtai and Yuri Gurevich.
\newblock Monotone versus positive.
\newblock {\em J. Assoc. Comput. Mach.}, 34(4):1004--1015, 1987.

\bibitem[AHM{\etalchar{+}}08]{AHMPS08}
Eric Allender, Lisa Hellerstein, Paul McCabe, Toniann Pitassi, and Michael~E.
  Saks.
\newblock Minimizing disjunctive normal form formulas and {AC}$^0$ circuits
  given a truth table.
\newblock {\em {SIAM} J. Comput.}, 38(1):63--84, 2008.

\bibitem[AK11]{akiyama_factors_2011}
Jin Akiyama and Mikio Kano.
\newblock {\em Factors and {Factorizations} of {Graphs}}, volume 2031 of {\em
  Lecture {Notes} in {Mathematics}}.
\newblock Springer Berlin Heidelberg, Berlin, Heidelberg, 2011.

\bibitem[AKR{\etalchar{+}}01]{AKRRV01}
Eric Allender, Michal Kouck{\'{y}}, Detlef Ronneburger, Sambuddha Roy, and
  V.~Vinay.
\newblock Time-space tradeoffs in the counting hierarchy.
\newblock In {\em Proceedings of the 16th Annual {IEEE} Conference on
  Computational Complexity, Chicago, Illinois, USA, June 18-21, 2001}, pages
  295--302. {IEEE} Computer Society, 2001.

\bibitem[AKS83]{aks83}
Mikl{\'{o}}s Ajtai, J{\'{a}}nos Koml{\'{o}}s, and Endre Szemer{\'{e}}di.
\newblock An {$O(n \log n)$} sorting network.
\newblock In {\em Proceedings of the 15th Annual {ACM} Symposium on Theory of
  Computing, 25-27 April, 1983, Boston, Massachusetts, {USA}}, pages 1--9.
  {ACM}, 1983.

\bibitem[And85]{andreev1985method}
Alexander~E Andreev.
\newblock On a method for obtaining lower bounds for the complexity of
  individual monotone functions.
\newblock {\em Doklady Akademii Nauk SSSR}, 282:1033--1037, 1985.

\bibitem[BCRV03]{playing_one}
Elmar Böhler, Nadia Creignou, Steffen Reith, and Heribert Vollmer.
\newblock Playing with boolean blocks, part i: Post’s lattice with
  applications to complexity theory.
\newblock {\em SIGACT News}, 2003.

\bibitem[BCRV04]{playing_two}
Elmar Böhler, Nadia Creignou, Steffen Reith, and Heribert Vollmer.
\newblock Playing with boolean blocks, part ii: Constraint satisfaction
  problems.
\newblock {\em ACM SIGACT-Newsletter}, 35, 2004.

\bibitem[BDHM92]{DBLP:journals/mst/BuntrockDHM92}
Gerhard Buntrock, Carsten Damm, Ulrich Hertrampf, and Christoph Meinel.
\newblock Structure and importance of logspace-mod class.
\newblock {\em Math. Syst. Theory}, 25(3):223--237, 1992.

\bibitem[Ber82]{berkowitz}
S.~J. Berkowitz.
\newblock On some relationships between monotone and non-monotone circuit
  complexity.
\newblock Technical Report. University of Toronto, 1982.

\bibitem[BGW99]{DBLP:journals/combinatorica/BabaiGW99}
L{\'{a}}szl{\'{o}} Babai, Anna G{\'{a}}l, and Avi Wigderson.
\newblock Superpolynomial lower bounds for monotone span programs.
\newblock {\em Combinatorica}, 19(3):301--319, 1999.

\bibitem[BHST14]{DBLP:conf/icalp/BlaisHST14}
Eric Blais, Johan H{\aa}stad, Rocco~A. Servedio, and Li{-}Yang Tan.
\newblock On {DNF} approximators for monotone boolean functions.
\newblock In {\em International Colloquium on Automata, Languages, and
  Programming \emph{(ICALP)}}, pages 235--246, 2014.

\bibitem[BL88]{DBLP:conf/crypto/Leichter88}
Josh~Cohen Benaloh and Jerry Leichter.
\newblock Generalized secret sharing and monotone functions.
\newblock In {\em Advances in Cryptology \emph{(CRYPTO)}}, pages 27--35, 1988.

\bibitem[BST13]{BST13}
Eric Blais, Dominik Scheder, and Li-Yang Tan.
\newblock Ajtai-gurevich redux. {M}anuscript.
\newblock 2013.

\bibitem[BT96]{DBLP:journals/jacm/BshoutyT96}
Nader~H. Bshouty and Christino Tamon.
\newblock On the {F}ourier spectrum of monotone functions.
\newblock {\em J. {ACM}}, 43(4):747--770, 1996.

\bibitem[Bul17]{DBLP:conf/focs/Bulatov17}
Andrei~A. Bulatov.
\newblock A dichotomy theorem for nonuniform csps.
\newblock In Chris Umans, editor, {\em 58th {IEEE} Annual Symposium on
  Foundations of Computer Science, {FOCS} 2017, Berkeley, CA, USA, October
  15-17, 2017}, pages 319--330. {IEEE} Computer Society, 2017.

\bibitem[Bul18]{DBLP:journals/siglog/Bulatov18}
Andrei~A. Bulatov.
\newblock Constraint satisfaction problems: complexity and algorithms.
\newblock {\em {ACM} {SIGLOG} News}, 5(4):4--24, 2018.

\bibitem[CDM21]{cdm21}
Arkadev Chattopadhyay, Rajit Datta, and Partha Mukhopadhyay.
\newblock Lower bounds for monotone arithmetic circuits via communication
  complexity.
\newblock In {\em {STOC} '21: 53rd Annual {ACM} {SIGACT} Symposium on Theory of
  Computing, Virtual Event, Italy, June 21-25, 2021}, pages 786--799. {ACM},
  2021.

\bibitem[CHO{\etalchar{+}}20]{choprs20}
Lijie Chen, Shuichi Hirahara, Igor~C. Oliveira, J{\'a}n Pich, Ninad Rajgopal,
  and Rahul Santhanam.
\newblock {Beyond Natural Proofs: Hardness Magnification and Locality}.
\newblock In {\em 11th Innovations in Theoretical Computer Science Conference
  (ITCS 2020)}, volume 151, pages 70:1--70:48, 2020.

\bibitem[CKR20]{CKR20}
Bruno~Pasqualotto Cavalar, Mrinal Kumar, and Benjamin Rossman.
\newblock Monotone circuit lower bounds from robust sunflowers.
\newblock In {\em {LATIN} 2020: Theoretical Informatics - 14th Latin American
  Symposium, S{\~{a}}o Paulo, Brazil, January 5-8, 2021, Proceedings}, volume
  12118 of {\em Lecture Notes in Computer Science}, pages 311--322. Springer,
  2020.

\bibitem[CKS01]{DBLP:books/daglib/0004131}
Nadia Creignou, Sanjeev Khanna, and Madhu Sudan.
\newblock {\em Complexity classifications of Boolean constraint satisfaction
  problems}, volume~7 of {\em {SIAM} monographs on discrete mathematics and
  applications}.
\newblock {SIAM}, 2001.

\bibitem[COS17]{cos15}
Xi~Chen, Igor~C. Oliveira, and Rocco~A. Servedio.
\newblock Addition is exponentially harder than counting for shallow monotone
  circuits.
\newblock In {\em S{TOC}'17---{P}roceedings of the 49th {A}nnual {ACM} {SIGACT}
  {S}ymposium on {T}heory of {C}omputing}, pages 1232--1245. ACM, New York,
  2017.

\bibitem[CZ16]{DBLP:conf/stoc/ChattopadhyayZ16}
Eshan Chattopadhyay and David Zuckerman.
\newblock Explicit two-source extractors and resilient functions.
\newblock In {\em Symposium on Theory of Computing \emph{(STOC)}}, pages
  670--683, 2016.

\bibitem[dRGR22]{DBLP:journals/sigact/RezendeGR22}
Susanna~F. de~Rezende, Mika G{\"{o}}{\"{o}}s, and Robert Robere.
\newblock Guest column: Proofs, circuits, and communication.
\newblock {\em {SIGACT} News}, 53(1):59--82, 2022.

\bibitem[FV98]{DBLP:journals/siamcomp/FederV98}
Tom{\'{a}}s Feder and Moshe~Y. Vardi.
\newblock The computational structure of monotone monadic {SNP} and constraint
  satisfaction: {A} study through datalog and group theory.
\newblock {\em {SIAM} J. Comput.}, 28(1):57--104, 1998.

\bibitem[GGLR98]{DBLP:conf/focs/GoldreichGLR98}
Oded Goldreich, Shafi Goldwasser, Eric Lehman, and Dana Ron.
\newblock Testing monotonicity.
\newblock In {\em Symposium on Foundations of Computer Science, \emph{(FOCS)}},
  pages 426--435, 1998.

\bibitem[GI12]{DBLP:journals/toc/GoldreichI12}
Oded Goldreich and Rani Izsak.
\newblock Monotone circuits: One-way functions versus pseudorandom generators.
\newblock {\em Theory Comput.}, 8(1):231--238, 2012.

\bibitem[GJW18]{gjw18}
Mika G{\"{o}}{\"{o}}s, Rahul Jain, and Thomas Watson.
\newblock Extension complexity of independent set polytopes.
\newblock {\em {SIAM} J. Comput.}, 47(1):241--269, 2018.

\bibitem[GKRS19]{gkrs19}
Mika G\"{o}\"{o}s, Pritish Kamath, Robert Robere, and Dmitry Sokolov.
\newblock Adventures in monotone complexity and {TFNP}.
\newblock In {\em 10th {I}nnovations in {T}heoretical {C}omputer {S}cience},
  volume 124 of {\em LIPIcs. Leibniz Int. Proc. Inform.}, pages Art. No. 38,
  19. Schloss Dagstuhl. Leibniz-Zent. Inform., Wadern, 2019.

\bibitem[GMOR15]{DBLP:conf/tcc/GuoMOR15}
Siyao Guo, Tal Malkin, Igor~C. Oliveira, and Alon Rosen.
\newblock The power of negations in cryptography.
\newblock In {\em Theory of Cryptography Conference \emph{(TCC)}}, pages
  36--65, 2015.

\bibitem[GS92]{gs92}
Michelangelo Grigni and Michael Sipser.
\newblock Monotone complexity.
\newblock In {\em Proceedings of the London Mathematical Society Symposium on
  Boolean Function Complexity}, page 57–75, USA, 1992. Cambridge University
  Press.

\bibitem[H{\aa}s86]{DBLP:conf/stoc/Hastad86}
Johan H{\aa}stad.
\newblock Almost optimal lower bounds for small depth circuits.
\newblock In {\em Proceedings of the 18th Annual {ACM} Symposium on Theory of
  Computing, May 28-30, 1986, Berkeley, California, {USA}}, pages 6--20, 1986.

\bibitem[Hir22]{Hir22}
Shuichi Hirahara.
\newblock {NP}-hardness of learning programs and partial {MCSP}.
\newblock In {\em Symposium on Foundations of Computer Science \emph{(FOCS)}},
  2022.

\bibitem[HWWY94]{hwwy94}
Johan H{\aa}stad, Ingo Wegener, Norbert Wurm, and Sang{-}Zin Yi.
\newblock Optimal depth, very small size circuits for symmetric functions in
  {AC}$^0$.
\newblock {\em Inf. Comput.}, 108(2):200--211, 1994.

\bibitem[IKL{\etalchar{+}}19]{DBLP:journals/jacm/IkenmeyerKLLMS19}
Christian Ikenmeyer, Balagopal Komarath, Christoph Lenzen, Vladimir Lysikov,
  Andrey Mokhov, and Karteek Sreenivasaiah.
\newblock On the complexity of hazard-free circuits.
\newblock {\em J. {ACM}}, 66(4):25:1--25:20, 2019.

\bibitem[JCG97]{10.1145/263867.263489}
Peter Jeavons, David Cohen, and Marc Gyssens.
\newblock Closure properties of constraints.
\newblock {\em J. ACM}, 44(4):527–548, jul 1997.

\bibitem[Jea98]{DBLP:journals/tcs/Jeavons98}
Peter Jeavons.
\newblock On the algebraic structure of combinatorial problems.
\newblock {\em Theor. Comput. Sci.}, 200(1-2):185--204, 1998.

\bibitem[JLL76]{DBLP:journals/mst/JonesLL76}
Neil~D. Jones, Y.~Edmund Lien, and William~T. Laaser.
\newblock New problems complete for nondeterministic log space.
\newblock {\em Math. Syst. Theory}, 10:1--17, 1976.

\bibitem[Joh03]{joh03}
Jan Johannsen.
\newblock The complexity of satisfiability problems with two occurrences.
\newblock Unpublished Manuscript, 2003.

\bibitem[Juk12]{jukna_2012}
Stasys Jukna.
\newblock {\em Boolean Function Complexity - Advances and Frontiers}.
\newblock Springer, 2012.

\bibitem[Kra97]{Krajicek97}
Jan Kraj\'{\i}\v{c}ek.
\newblock Interpolation theorems, lower bounds for proof systems, and
  independence results for bounded arithmetic.
\newblock {\em J. Symbolic Logic}, 62(2):457--486, 1997.

\bibitem[KS92]{DBLP:journals/siamdm/KalyanasundaramS92}
Bala Kalyanasundaram and Georg Schnitger.
\newblock The probabilistic communication complexity of set intersection.
\newblock {\em {SIAM} J. Discret. Math.}, 5(4):545--557, 1992.

\bibitem[Kup21]{DBLP:conf/lics/Kuperberg21}
Denis Kuperberg.
\newblock Positive first-order logic on words.
\newblock In {\em Symposium on Logic in Computer Science \emph{(LICS)}}, pages
  1--13, 2021.

\bibitem[Kup22]{DBLP:journals/corr/abs-2201-11619}
Denis Kuperberg.
\newblock Positive first-order logic on words and graphs.
\newblock {\em CoRR}, abs/2201.11619, 2022.

\bibitem[KW93]{DBLP:conf/coco/KarchmerW93}
Mauricio Karchmer and Avi Wigderson.
\newblock On span programs.
\newblock In {\em Structure in Complexity Theory Conference \emph{(CCC)}},
  pages 102--111, 1993.

\bibitem[LT09]{DBLP:journals/tcs/LaroseT09}
Beno{\^{\i}}t Larose and Pascal Tesson.
\newblock Universal algebra and hardness results for constraint satisfaction
  problems.
\newblock {\em Theor. Comput. Sci.}, 410(18):1629--1647, 2009.

\bibitem[Oko82]{okolnishnikova_1982}
E.~A. Okol'nishnikova.
\newblock The effect of negations on the complexity of realization of monotone
  {B}oolean functions by formulas of bounded depth.
\newblock {\em Metody Diskret. Analiz.}, (38):74--80, 1982.

\bibitem[Oli15]{Oliveira15}
Igor~C. Oliveira.
\newblock Unconditional lower bounds in complexity theory, 2015.
\newblock (PhD Thesis, Columbia University).

\bibitem[OSS19]{DBLP:conf/coco/OliveiraS019}
Igor~C. Oliveira, Rahul Santhanam, and Srikanth Srinivasan.
\newblock Parity helps to compute majority.
\newblock In {\em 34th Computational Complexity Conference, {CCC} 2019, July
  18-20, 2019, New Brunswick, NJ, {USA}}, pages 23:1--23:17, 2019.

\bibitem[Pos41]{post41}
Emil~L. Post.
\newblock {\em The Two-Valued Iterative Systems of Mathematical Logic. (AM-5)}.
\newblock Princeton University Press, 1941.

\bibitem[PR17]{DBLP:conf/stoc/PitassiR17}
Toniann Pitassi and Robert Robere.
\newblock Strongly exponential lower bounds for monotone computation.
\newblock In {\em Symposium on Theory of Computing \emph{(STOC)}}, pages
  1246--1255, 2017.

\bibitem[Pud97]{DBLP:journals/jsyml/Pudlak97}
Pavel Pudl{\'{a}}k.
\newblock Lower bounds for resolution and cutting plane proofs and monotone
  computations.
\newblock {\em J. Symb. Log.}, 62(3):981--998, 1997.

\bibitem[Qui53]{quine_1953}
W.~V. Quine.
\newblock Two theorems about truth functions.
\newblock {\em Bol. Soc. Mat. Mexicana}, 10:64--70, 1953.

\bibitem[Raz85a]{razborov_boolean_85}
A.~A. Razborov.
\newblock Lower bounds on the monotone complexity of some {B}oolean functions.
\newblock {\em Dokl. Akad. Nauk SSSR}, 281(4):798--801, 1985.

\bibitem[Raz85b]{razborov1985lower}
Alexander~A Razborov.
\newblock Lower bounds on monotone complexity of the logical permanent.
\newblock {\em Mathematical Notes of the Academy of Sciences of the USSR},
  37(6):485--493, 1985.

\bibitem[Rei05]{DBLP:conf/stoc/Reingold05}
Omer Reingold.
\newblock Undirected st-connectivity in log-space.
\newblock In Harold~N. Gabow and Ronald Fagin, editors, {\em Proceedings of the
  37th Annual {ACM} Symposium on Theory of Computing, Baltimore, MD, USA, May
  22-24, 2005}, pages 376--385. {ACM}, 2005.

\bibitem[Rez23]{susanna_2023}
Susanna Rezende.
\newblock Personal communication, 2023.

\bibitem[RM99]{DBLP:journals/combinatorica/RazM99}
Ran Raz and Pierre McKenzie.
\newblock Separation of the monotone {NC} hierarchy.
\newblock {\em Combinatorica}, 19(3):403--435, 1999.

\bibitem[Ros08a]{rossman_2008}
Benjamin Rossman.
\newblock Homomorphism preservation theorems.
\newblock {\em J. ACM}, 55(3):Art. 15, 53, 2008.

\bibitem[Ros08b]{DBLP:conf/stoc/Rossman08}
Benjamin Rossman.
\newblock On the constant-depth complexity of $k$-clique.
\newblock In {\em Symposium on Theory of Computing \emph{(STOC)}}, pages
  721--730, 2008.

\bibitem[Ros17a]{rossman_entropy}
Benjamin Rossman.
\newblock An entropy proof of the switching lemma and tight bounds on the
  decision-tree size of {AC}$^0$.
\newblock 2017.

\bibitem[Ros17b]{rossman_16}
Benjamin Rossman.
\newblock An improved homomorphism preservation theorem from lower bounds in
  circuit complexity.
\newblock In {\em Innovations in Theoretical Computer Science Conference
  \emph{(ITCS)}}, pages 27:1--27:17, 2017.

\bibitem[RPRC16]{rprc16}
Robert Robere, Toniann Pitassi, Benjamin Rossman, and Stephen~A. Cook.
\newblock Exponential lower bounds for monotone span programs.
\newblock In {\em {IEEE} 57th Annual Symposium on Foundations of Computer
  Science, {FOCS} 2016, 9-11 October 2016, Hyatt Regency, New Brunswick, New
  Jersey, {USA}}, pages 406--415. {IEEE} Computer Society, 2016.

\bibitem[RW92]{RW92}
Ran Raz and Avi Wigderson.
\newblock Monotone circuits for matching require linear depth.
\newblock {\em J. {ACM}}, 39(3):736--744, 1992.

\bibitem[Sch78]{DBLP:conf/stoc/Schaefer78}
Thomas~J. Schaefer.
\newblock The complexity of satisfiability problems.
\newblock In {\em Proceedings of the 10th Annual {ACM} Symposium on Theory of
  Computing, May 1-3, 1978, San Diego, California, {USA}}, pages 216--226.
  {ACM}, 1978.

\bibitem[Sto95]{DBLP:conf/lics/Stolboushkin95}
Alexei~P. Stolboushkin.
\newblock Finitely monotone properties.
\newblock In {\em Symposium on Logic in Computer Science \emph{(LICS)}}, pages
  324--330, 1995.

\bibitem[Tar88]{tardos_1988}
\'{E}. Tardos.
\newblock The gap between monotone and nonmonotone circuit complexity is
  exponential.
\newblock {\em Combinatorica}, 8(1):141--142, 1988.

\bibitem[Urq87]{DBLP:journals/jacm/Urquhart87}
Alasdair Urquhart.
\newblock Hard examples for resolution.
\newblock {\em J. {ACM}}, 34(1):209--219, 1987.

\bibitem[Val80]{DBLP:journals/tcs/Valiant80}
Leslie~G. Valiant.
\newblock Negation can be exponentially powerful.
\newblock {\em Theor. Comput. Sci.}, 12:303--314, 1980.

\bibitem[Zhu17]{DBLP:conf/focs/Zhuk17}
Dmitriy Zhuk.
\newblock A proof of {CSP} dichotomy conjecture.
\newblock In Chris Umans, editor, {\em 58th {IEEE} Annual Symposium on
  Foundations of Computer Science, {FOCS} 2017, Berkeley, CA, USA, October
  15-17, 2017}, pages 331--342. {IEEE} Computer Society, 2017.

\end{thebibliography}
\end{document}
